\documentclass[11pt,a4paper]{article}

\usepackage[T1]{fontenc}
\usepackage[utf8]{inputenc}
\usepackage{lmodern}
\usepackage[margin=2.6cm]{geometry}
\usepackage{amsmath,amssymb,amsthm}
\usepackage{mathtools}
\usepackage{booktabs}
\usepackage{array}
\usepackage{longtable}
\usepackage{microtype}
\usepackage{enumitem}
\usepackage{alltt}
\usepackage{newunicodechar}
\usepackage[hidelinks]{hyperref}

\allowdisplaybreaks
\sloppy
\emergencystretch=3em

% Unicode characters occurring in the reproduced Lean source.
\newunicodechar{β}{\ensuremath{\beta}}
\newunicodechar{η}{\ensuremath{\eta}}
\newunicodechar{μ}{\ensuremath{\mu}}
\newunicodechar{ℂ}{\ensuremath{\mathbb{C}}}
\newunicodechar{ℕ}{\ensuremath{\mathbb{N}}}
\newunicodechar{ℝ}{\ensuremath{\mathbb{R}}}
\newunicodechar{→}{\ensuremath{\to}}
\newunicodechar{↦}{\ensuremath{\mapsto}}
\newunicodechar{∂}{\ensuremath{\partial}}
\newunicodechar{∑}{\ensuremath{\textstyle\sum}}
\newunicodechar{∞}{\ensuremath{\infty}}
\newunicodechar{∫}{\ensuremath{\textstyle\int}}
\newunicodechar{≤}{\ensuremath{\le}}
\newunicodechar{≥}{\ensuremath{\ge}}
\newunicodechar{⨆}{\ensuremath{\textstyle\bigsqcup}}

% A Lean code block: verbatim-like, but macros and Unicode still work.
\newenvironment{leancode}{\par\smallskip\begin{alltt}\small}{\end{alltt}\smallskip}

\newtheorem{theorem}{Theorem}[section]
\newtheorem{proposition}[theorem]{Proposition}
\newtheorem{lemma}[theorem]{Lemma}
\newtheorem{corollary}[theorem]{Corollary}
\theoremstyle{definition}
\newtheorem{definition}[theorem]{Definition}
\newtheorem{remark}[theorem]{Remark}
\newtheorem{convention}[theorem]{Convention}

\DeclareMathOperator{\Var}{Var}
\DeclareMathOperator*{\esssup}{ess\,sup}
\newcommand{\RR}{\mathbb{R}}
\newcommand{\CC}{\mathbb{C}}
\newcommand{\NN}{\mathbb{N}}
\newcommand{\QQ}{\mathbb{Q}}
\newcommand{\EE}{\mathbb{E}}
\newcommand{\PP}{\mathbb{P}}
\newcommand{\ii}{\mathrm{i}}
\newcommand{\dd}{\,\mathrm{d}}
\newcommand{\kap}{\kappa}
\newcommand{\Lean}[1]{\texttt{#1}}
\newcommand{\file}[1]{\texttt{#1}}

% A compact "formalized as" marker.
\newcommand{\formal}[2]{%
  \par\smallskip
  \begingroup\small\raggedright\emergencystretch=5em
  \noindent\textit{In Lean:} \Lean{#1} in \file{#2}.\par
  \endgroup\smallskip}

\title{The i.i.d.\ Berry--Esseen constant is at most $0.423$:\\
a machine-checked proof}

\author{}
\date{13 September 2026}

\begin{document}

\maketitle

\begin{abstract}
Let $C_{\mathrm{BE}}$ be the least constant such that
$\sup_x|F_n(x)-\Phi(x)|\le C_{\mathrm{BE}}\,\beta_3/\sqrt n$ for every centered,
variance-one law with finite third absolute moment $\beta_3$ and every sample
size $n\ge 1$, where $F_n$ is the distribution function of the normalized sum of
$n$ independent copies and $\Phi$ is the standard normal distribution function.
We describe a complete proof of $C_{\mathrm{BE}}\le 0.423$ which has been
formalized in Lean~4 on top of Mathlib, in the module tree
\file{BerryEsseen/Upper0423/} of the accompanying repository, and whose final
statement is
\begin{center}
  \small
  \Lean{BerryEsseen.Upper0423.berryEsseenConstant\_le\_0423}\\[2pt]
  ${}:{}$\Lean{BerryEsseen.berryEsseenConstant} $\le$ \Lean{0.423}
\end{center}
The proof combines the Prawitz smoothing inequality with a family of pointwise
characteristic-function estimates in the three moment coordinates
$\beta=\EE|X|^3$, $\eta=\tfrac12\EE|X-X'|^3-\beta$ and $s=\EE|X|$, an
interval/Taylor-jet checking engine whose soundness is proved inside Lean, and
two families of finite certificates: $569$ per-sample-size parameter covers with
$18\,845$ leaves for $2\le n\le 255$, and a single certificate of $281$ cells
that is uniform in $n$ for all $n\ge 256$. The regimes $n=1$ and
$\beta/\sqrt n\ge 1.3003$ are handled analytically with no computation. There is
no \Lean{sorry} and no user-declared axiom; the finite certificates are
discharged by \Lean{native\_decide}, so they additionally trust Lean's compiler.
We state the trust boundary precisely and identify which analytic ingredients
originate in the earlier work of Xiao and Li.
\end{abstract}

\tableofcontents

\bigskip

\section{Introduction}
\label{sec:intro}

\subsection{The constant}

Let $X,X_1,X_2,\dots$ be independent real random variables with a common law
$\mu$ satisfying
\begin{equation}
  \label{eq:admissible}
  \EE X = 0,\qquad \EE X^2 = 1,\qquad \beta:=\EE|X|^3<\infty ,
\end{equation}
and set
\[
  S_n=\frac{X_1+\dots+X_n}{\sqrt n},\qquad
  F_n(x)=\PP(S_n\le x),\qquad
  \Delta_n=\sup_{x\in\RR}\bigl|F_n(x)-\Phi(x)\bigr| .
\]
The Berry--Esseen theorem \cite{Berry1941,Esseen1942} asserts the existence of a
universal constant $C$ with $\Delta_n\le C\beta/\sqrt n$ for all $n\ge1$ and all
laws satisfying \eqref{eq:admissible}. We write $C_{\mathrm{BE}}$ for the least
such constant, i.e.
\begin{equation}
  \label{eq:constant}
  C_{\mathrm{BE}}
  =\sup\Bigl\{\tfrac{\sqrt n}{\beta}\,\bigl|F_n(x)-\Phi(x)\bigr|
  \;:\;\mu \text{ satisfies } \eqref{eq:admissible},\ n\ge1,\ x\in\RR\Bigr\}
  \in[0,\infty] .
\end{equation}
Esseen \cite{Esseen1956} proved the lower bound
\[
  C_{\mathrm{BE}}\;\ge\;c_{\mathrm E}:=\frac{\sqrt{10}+3}{6\sqrt{2\pi}}
  =0.4097321\ldots,
\]
attained asymptotically along suitably centered and normalized Bernoulli sums,
and it is conjectured (usually attributed to Zolotarev \cite{Zolotarev1967}) that
$C_{\mathrm{BE}}=c_{\mathrm E}$. In the special case of Bernoulli (equivalently
binomial) summands the value $c_{\mathrm E}$ is known to be optimal; see
Schulz \cite{Schulz2016} and Zolotukhin, Nagaev and Chebotarev
\cite{Zolotukhin2018}.

Table~\ref{tab:history} lists the published upper bounds. The bibliographic data
follow the survey entry \cite{Tao2026}.

\begin{table}[htbp]
  \centering
  \begin{tabular}{llp{7.4cm}}
    \toprule
    Bound & Reference & Comment\\
    \midrule
    $0.82$   & Zolotarev \cite{Zolotarev1967}        & Zolotarev-type smoothing inequalities; also $0.9051$ for independent (non-i.i.d.) summands.\\
    $0.7975$ & van Beek \cite{vanBeek1972}            & Fourier-analytic refinement.\\
    $0.7655$ & Shiganov \cite{Shiganov1986}           & \\
    $0.7056$ & Shevtsova \cite{Shevtsova2006}         & \\
    $0.5129$ & Korolev--Shevtsova \cite{KorolevShevtsova2009} & From the structural bound $\Delta_n\le0.34445(\beta+0.489)/\sqrt n$.\\
    $0.4785$ & Tyurin \cite{Tyurin2009}               & \\
    $0.4748$ & Shevtsova \cite{Shevtsova2011}         & \\
    $0.4690$ & Shevtsova \cite{Shevtsova2013}         & From
      $\Delta_n\le\min\{0.4690\,\beta,\ 0.3322(\beta+0.429),\ 0.3031(\beta+0.646)\}/\sqrt n$.\\
    \bottomrule
  \end{tabular}
  \caption{Published upper bounds for $C_{\mathrm{BE}}$.}
  \label{tab:history}
\end{table}

\subsection{The result}

\begin{theorem}
  \label{thm:main}
  $C_{\mathrm{BE}}\le \dfrac{423}{1000}$.
\end{theorem}

Theorem~\ref{thm:main} is formalized in Lean~4 over Mathlib \cite{Mathlib}. The
constant \eqref{eq:constant} is defined in the repository as an extended
nonnegative real supremum, with finiteness a consequence rather than a hypothesis:

\begin{leancode}
def berryEsseenConstant : ℝ≥0∞ :=
  ⨆ (μ : ProbabilityMeasure ℝ) (_ : Admissible μ) (n : ℕ) (_ : 0 < n) (x : ℝ),
    normalizedError μ n x
\end{leancode}

\noindent and the theorem proved is
\begin{leancode}
theorem BerryEsseen.Upper0423.berryEsseenConstant_le_0423 :
    BerryEsseen.berryEsseenConstant ≤ 0.423
\end{leancode}
in \file{BerryEsseen/Upper0423/Main.lean}. Section~\ref{sec:notation} reproduces
all the ambient definitions verbatim so that the statement can be audited without
reading Lean beyond these few lines.

\subsection{What is proved and what is computed}

The proof has two very different kinds of ingredient, and we keep them
rigorously separated throughout.

\begin{enumerate}[label=(\alph*),leftmargin=2.2em]
\item \emph{Analysis.} The Prawitz smoothing inequality, a list of pointwise
  estimates for the characteristic function of an admissible law in terms of
  $(\beta,\eta,s)$, the geometry of the moment region, elementary Gaussian
  bounds, the semantics of a fixed-point interval arithmetic and of fourth-order
  Taylor jets, the soundness of a Taylor/range quadrature rule, and the algebra
  that turns a certified numerical budget into the final inequality. All of this
  is proved in Lean by ordinary means and is checked by the Lean kernel.

\item \emph{Finite computation.} Exactly $579$ closed Boolean statements of the
  form ``this explicit certificate is accepted by this explicit checker'' are
  discharged by the tactic \Lean{native\_decide}, which compiles the Boolean
  expression to machine code and runs it. Each such use introduces one auxiliary
  axiom, so the finite part of the proof trusts the Lean compiler and runtime in
  addition to the kernel.
\end{enumerate}

Nothing about the \emph{search} for the certificates is trusted: the search was
done outside Lean, but every accepted numerical bound is recomputed by a checker
whose soundness is a Lean theorem, quantified over \emph{all} inputs the checker
accepts. Section~\ref{sec:trust} states the trust boundary precisely and lists
what we do and do not claim.

\subsection{Outline}

Section~\ref{sec:attribution} records the provenance of the analytic
ingredients. Section~\ref{sec:notation} fixes notation and reproduces the formal
definitions. Section~\ref{sec:moments} develops the moment geometry
($\beta\ge1$, $0\le\eta\le s\le1$, $s\beta\ge1$, $\beta+\eta\ge2$).
Section~\ref{sec:prawitz} states the smoothing inequality as formalized.
Section~\ref{sec:chf} collects the characteristic-function estimates.
Section~\ref{sec:split} gives the case decomposition, and
Sections~\ref{sec:one} and \ref{sec:largeell} dispose of the two elementary
cases. Section~\ref{sec:engine} describes the interval/jet engine and its
soundness, Section~\ref{sec:covers} the per-sample-size covers for
$2\le n\le255$, and Section~\ref{sec:uniform} the uniform certificate for
$n\ge256$. Section~\ref{sec:trust} is the trust boundary and
Section~\ref{sec:repro} the reproduction instructions.
Section~\ref{sec:related} discusses related and concurrent work.

\section{Provenance and attribution}
\label{sec:attribution}

Because this is a formalization of a computer-assisted argument that builds on
earlier work, it is important to say at the outset which mathematical
ingredients are inherited and which are not.

\subsection{Ingredients originating in Xiao--Li}

A substantial part of the analytic framework used here originates in the work of
Haonan Xiao and Chunlin Li \cite{XiaoLi2026}, specifically the repository
\begin{center}
\url{https://github.com/haonan-xiao/iid-berry-esseen}
\end{center}
at commit \texttt{f00d781c8c42908bf4aef2739129442a0919fa38}, which proves the
bound $C_{\mathrm{BE}}\le 879/2000=0.4395$ with its own Lean formalization
(accessed 13 September 2026). We refer to that work as [XL]. The following
ingredients are due to [XL] and are used here:

\begin{enumerate}[label=(\roman*),leftmargin=2.4em]
\item \textbf{The moment coordinates and the moment region.} The idea of
  carrying the symmetrized third-moment ratio alongside $\beta$ and the first
  absolute moment $s=\EE|X|$, and of using the resulting region as the parameter
  domain of the whole argument, is from [XL], \file{proof-guide.md}
  \S\S2--3. Our $\eta$ is an affine reparametrization of their ratio
  $r=\EE|X-X'|^3/(2\EE|X|^3)$, namely $\eta=\beta(r-1)$.

\item \textbf{The Prawitz normalization convention.} The particular form of the
  doubled Prawitz kernel and of the resulting exact upper functional
  (Section~\ref{sec:prawitz}) follows [XL], \file{proof-guide.md} \S\S5--6.

\item \textbf{The stop-loss identity and the excess bounds.} The identity
  expressing $\EE|X-X'|^3$ through the stop-loss transform, and the resulting
  inequalities $\eta\ge0$ and $\beta+\eta\ge2$, are from [XL],
  \file{lean/BerryEsseen/Moments/ThirdMomentRatio.lean}
  (\Lean{symmetrizedThirdAbsoluteMoment\_stopLoss\_identity}). See
  Section~\ref{sec:moments}.

\item \textbf{The two-point sine circle.} The mechanism by which a centered law
  is decomposed into centered two-point laws and the imaginary part of the
  characteristic function is confined to a circle jointly with an energy
  coordinate is from [XL],
  \file{lean/BerryEsseen/CharacteristicFunctions/SineRemainder.lean}
  (\Lean{sine\_remainder\_circle} together with the two-point energy and moment
  lemmas). See Proposition~\ref{prop:circle}.

\item \textbf{The radial estimates.} The separate radial bounds for the real and
  imaginary parts in terms of $1-s$ are from [XL],
  \file{lean/BerryEsseen/CharacteristicFunctions/TwoPointComparison.lean}
  (\Lean{charFun\_real\_radial\_bound}, \Lean{charFun\_imag\_radial\_bound}).
  See Proposition~\ref{prop:radial}.

\item \textbf{The disk/strip intersection.} The idea of intersecting the
  second-order Taylor disk with the real and imaginary strips, which is what the
  \emph{coordinate} modulus branch of our engine implements, is from [XL],
  \file{lean/BerryEsseen/CharacteristicFunctions/MomentDisk.lean}.

\item \textbf{The cubic cosine constant and the convex profile.} The constant
  \[
    \kap=\max_{v>0}\frac{\cos v-1+v^2/2}{v^3}
  \]
  and its use as a global majorant appear in [XL], whose proof guide brackets it
  as $0.09916191350<\kap<0.09916191353$ and uses the rational upper bound
  $\Lean{kappaPlus}=9916191353/10^{11}$. The convex profile obtained by splicing
  the supporting line $\tfrac12-\kap v$ to $(1-\cos v)/v^2$ also appears
  there. Our development uses the coarser rational majorant
  $\kap\le 49581/500000=0.099162$, which is what the certificates are calibrated
  to. See Section~\ref{sec:kappa}.
\end{enumerate}

Our Lean proofs of all these lemmas were written from scratch in the present
development and are checked by the Lean kernel; we do not import, translate or
otherwise reuse [XL]'s Lean code. What is inherited is mathematical: the choice
of coordinates, the smoothing normalization, and the specific mechanisms listed
above. We make no claim of priority for any of them.

\subsection{What is not used}

Theorem~\ref{thm:main} does not use [XL]'s numerical result, nor any of their
certificates, nor their region decomposition. This is not a stylistic choice:
their certificates are budgets for the target $0.4395$ and cannot be reused as
budgets for $0.423$, and their Lyapunov cutoff $\beta/\sqrt n\le 56/45$ is
tailored to $0.4395$ (our analytic cap $\Delta_n\le 11/20$ forces the smaller
cutoff $13003/10000$; see Section~\ref{sec:largeell}). We make no assertion about
the correctness of their $0.4395$ theorem, which we have not audited.

Theorem~\ref{thm:main} also does \emph{not} use Shevtsova's published bound
\cite{Shevtsova2011,Shevtsova2013}, or any other result from
Table~\ref{tab:history}. Earlier drafts of the underlying research notes did
invoke Shevtsova's estimate to dispose of the large-third-moment region; the
formalized proof replaces that input by the direct treatment of
Sections~\ref{sec:largeell} and \ref{sec:covers}. The only external mathematical
library used is Mathlib.

\subsection{What is new here}

Relative to [XL] the following ingredients are, as far as we can tell, new in the
present development. We state this descriptively and claim no more.

\begin{enumerate}[label=(\alph*),leftmargin=2.2em]
\item \textbf{The anchored imaginary envelope.} The cubic imaginary bound
  $|\operatorname{Im}f(u)|\le\tfrac16\sqrt{(\beta-1)(\beta+5/3)}\,|u|^3$
  anchored at $\cos u$ rather than at~$1$
  (Proposition~\ref{prop:imcubic}), and the sharpened quadratic bound
  $|\operatorname{Im}f(u)|\le\tfrac12\sqrt{2(\beta-1)/3}\,u^2$
  (Proposition~\ref{prop:imquad}).
\item \textbf{The signed real interval.} The two-sided estimate
  $\cos u\le\operatorname{Re}f(u)$ on $[0,\pi]$
  (Proposition~\ref{prop:relower}) together with the three upper bounds of
  Proposition~\ref{prop:reupper}, which lets the engine use the
  \emph{signed} real part in the coordinate modulus branch.
\item \textbf{The joint complex anchor disk.} The ODE-based bound
  $\|f(u)-\cos u\|\le\sqrt{(\beta-1)(\beta+5/3)}\,(u-\sin u)$ on $[0,\pi]$
  (Proposition~\ref{prop:ode}).
\item \textbf{The geometric-sum comparison.} The replacement of the crude
  $|a^n-b^n|\le n|a-b|$ by the exact telescoping sum
  $\sum_{j<n}|a|^j|b|^{n-1-j}$, evaluated by a halving recursion inside the
  certified jet arithmetic (Section~\ref{sec:lowfactor}).
\item \textbf{The uniform-in-$n$ certificate.} A single finite certificate
  covering \emph{all} $n\ge256$ simultaneously, obtained by reducing the
  smoothing functional to five one-dimensional integrals in the two scalars
  $L=(\beta+\eta)/\sqrt n$ and $h=1/\sqrt n$ (Section~\ref{sec:uniform}).
\item \textbf{The elimination of the published external input}, as described
  above.
\end{enumerate}

\section{Notation and the formal statement}
\label{sec:notation}

\subsection{The ambient definitions}

The following definitions are those of \file{BerryEsseen/Defs.lean} (identical
to the independently auditable copy in \file{Challenge.lean}), reproduced here
verbatim.

\begin{leancode}
def thirdAbsMoment (μ : ProbabilityMeasure ℝ) : ℝ := ∫ x, |x| ^ 3 ∂(μ : Measure ℝ)

structure Admissible (μ : ProbabilityMeasure ℝ) : Prop where
  integrable_id : Integrable (fun x : ℝ ↦ x) (μ : Measure ℝ)
  integrable_sq : Integrable (fun x : ℝ ↦ x ^ 2) (μ : Measure ℝ)
  integrable_abs_cube : Integrable (fun x : ℝ ↦ |x| ^ 3) (μ : Measure ℝ)
  mean_zero : (∫ x, x ∂(μ : Measure ℝ)) = 0
  second_moment_one : (∫ x, x ^ 2 ∂(μ : Measure ℝ)) = 1

def normalizedSum (n : ℕ) (x : Fin n → ℝ) : ℝ := (∑ i, x i) / Real.sqrt n

def normalizedSumLaw (μ : ProbabilityMeasure ℝ) (n : ℕ) : Measure ℝ :=
  Measure.map (normalizedSum n) (Measure.pi fun _ : Fin n ↦ (μ : Measure ℝ))

def normalCDF (x : ℝ) : ℝ := (gaussianReal 0 1).real (Set.Iic x)

def normalizedError (μ : ProbabilityMeasure ℝ) (n : ℕ) (x : ℝ) : ℝ≥0∞ :=
  ENNReal.ofReal (Real.sqrt n / thirdAbsMoment μ *
    |(normalizedSumLaw μ n).real (Set.Iic x) - normalCDF x|)

def berryEsseenConstant : ℝ≥0∞ :=
  ⨆ (μ : ProbabilityMeasure ℝ) (_ : Admissible μ) (n : ℕ) (_ : 0 < n) (x : ℝ),
    normalizedError μ n x
\end{leancode}

Three points deserve emphasis. First, independence is realized by the finite
product measure $\mu^{\otimes n}$ on $\RR^{\mathrm{Fin}\,n}$, so no auxiliary
probability space is assumed to exist. Second, the supremum is taken in
$[0,\infty]$, so finiteness of $C_{\mathrm{BE}}$ is part of the conclusion and
not a hypothesis. Third, \Lean{normalizedError} is defined with
$\sqrt n/\beta$ as a multiplicative prefactor rather than by an inequality
$\Delta_n\le C\beta/\sqrt n$; the two readings agree because $\beta>0$ for every
admissible law (Lemma~\ref{lem:betapos}).

The reduction from the supremum to a pointwise statement is the trivial
\begin{equation}
  \label{eq:constant-le-iff}
  C_{\mathrm{BE}}\le C
  \iff
  \forall \mu\ \text{admissible},\ \forall n\ge1,\ \forall x,\quad
  \Lean{normalizedError}\ \mu\ n\ x\le C ,
\end{equation}
which is \Lean{BerryEsseen.constant\_le\_iff} in \file{BerryEsseen/Constant.lean}.
Everything below therefore concerns the single pointwise quantity
\[
  \mathcal E(\mu,n,x):=\frac{\sqrt n}{\beta}\bigl|F_n(x)-\Phi(x)\bigr| ,
\]
which we must bound by $423/1000$.

\subsection{Moment coordinates}

Throughout, $\mu$ is an admissible law, $X\sim\mu$, $X'$ an independent copy, and
$f(u)=\EE e^{\ii uX}$ is the characteristic function. We use the three moment
coordinates
\begin{equation}
  \label{eq:coords}
  \beta=\EE|X|^3,\qquad
  \eta=\tfrac12\EE|X-X'|^3-\beta,\qquad
  s=\EE|X| ,
\end{equation}
and the derived quantity
\begin{equation}
  \label{eq:gamma}
  \Gamma=\sqrt{(\beta-1)^2+\beta^2-\eta^2} .
\end{equation}
We call $\eta$ the \emph{symmetrization excess}; it is the excess of the
half-symmetrized third absolute moment over $\beta$ itself, and it vanishes for
the Rademacher law. Carrying $(\beta,\eta,s)$ jointly --- equivalently, carrying
$\beta$, the symmetrized third-moment ratio $r=1+\eta/\beta$ and the first
absolute moment --- is the choice of coordinates made by [XL] (item (i) of
Section~\ref{sec:attribution}), and the region \eqref{eq:region} below is their
moment region.
\formal{thirdAbsMoment, symmExcess, absMean, gammaR}{BerryEsseen/Defs.lean, Upper0423/Analytic/Moments.lean, Upper0423/LargeN/Defs.lean}

We also write
\[
  \ell=\frac{\beta}{\sqrt n}
\]
for the Lyapunov fraction, and reserve $g(u)=e^{-u^2/2}$ for the standard normal
characteristic function.

\subsection{Constants}

Table~\ref{tab:constants} lists all the numerical constants that occur in the
proof; each is a rational number written exactly as it appears in the Lean
sources.

\begin{table}[htbp]
  \centering
  \begin{tabular}{llp{7.6cm}}
    \toprule
    Symbol & Exact value & Role\\
    \midrule
    target & $423/1000$ & the bound being proved (\Lean{Native.target})\\
    $\kap$ & $49581/500000$ & rational majorant for $\max_v(\cos v-1+v^2/2)/v^3$ (\Lean{Native.kappa}, \Lean{kappaR})\\
    $B$ & $712569/454000$ & upper end of the original small-$\beta$ covers (\Lean{deltaMax}${}=B-1=258569/454000$)\\
    $\varepsilon$ & $1/1000$ & low-frequency omission cutoff (\Lean{Native.epsilon})\\
    --- & $13003/10000$ & Lyapunov threshold for the Cantelli case\\
    --- & $11/20$ & the universal Cantelli cap for $\Delta_n$\\
    --- & $21/50$ & the $n=1$ constant\\
    $A$ & $(8-64\kap)/(1-\cos4)$ & scale of the curved branch of the convex profile\\
    $2^{40}$ & $1099511627776$ & fixed-point scale of the interval arithmetic (\Lean{Native.scale})\\
    $\tau$ & $1/4$ & smoothing split in the uniform certificate\\
    $r$ & $16$ & $r^2=256$, the first sample size of the uniform regime\\
    $m$ & $256$ & panel count per cell in the uniform certificate\\
    $\nu$ & $\bigl(\tfrac{22/7}{3(1-(22/7)^2/96)}\bigr)^2\!/2$ & quadratic coefficient of the normal-correction majorant (\Lean{LargeN.nu})\\
    \bottomrule
  \end{tabular}
  \caption{The numerical constants, exactly as in the Lean sources.}
  \label{tab:constants}
\end{table}

\section{The moment region}
\label{sec:moments}

All the pointwise estimates of Section~\ref{sec:chf} are stated in terms of
$(\beta,\eta,s)$, so the first task is to delimit the region in which these
parameters live. The following four facts are what the certificates use.

\begin{lemma}[$\beta\ge1$]
  \label{lem:betapos}
  For every admissible $\mu$ we have $\beta\ge1$, and in particular $\beta>0$.
\end{lemma}

\begin{proof}
  The polynomial $w(y)=(y-1)^2(2y+1)=2y^3-3y^2+1$ is nonnegative for $y\ge0$.
  Applying it to $y=|X|$ and using $\EE X^2=1$ gives
  $0\le\EE w(|X|)=2\beta-3+1=2(\beta-1)$.
\end{proof}
\formal{momentExcessWeight, one\_le\_thirdAbsMoment, thirdAbsMoment\_pos}{BerryEsseen/MomentExcess.lean}

\begin{lemma}[$0\le\eta\le s\le1$ and $s\beta\ge1$]
  \label{lem:eta}
  For every admissible $\mu$,
  \[
    0\le\eta\le s\le 1,\qquad s\beta\ge1 .
  \]
\end{lemma}

\begin{proof}
  For $s\le1$ apply Cauchy--Schwarz: $s=\EE|X|\le(\EE X^2)^{1/2}=1$. For
  $s\beta\ge1$ apply Cauchy--Schwarz in the form
  $1=(\EE X^2)^2=(\EE|X|^{1/2}\cdot|X|^{3/2})^2\le\EE|X|\cdot\EE|X|^3=s\beta$.
  For $\eta\ge0$ and $\eta\le s$ we use two elementary scalar inequalities:
  \[
    |x-y|^3\ge|x|^3+|y|^3-3xy(|x|+|y|),
    \qquad
    |x-y|^3\le(x-y)^2(|x|+|y|) .
  \]
  Integrating the first against $\mu\otimes\mu$ and using $\EE X=0$ gives
  $\EE|X-X'|^3\ge2\beta$, i.e. $\eta\ge0$; integrating the second and expanding
  $(x-y)^2$ gives $\EE|X-X'|^3\le2\beta+2s$, i.e. $\eta\le s$.
\end{proof}
\formal{abs\_sub\_cube\_ge, abs\_sub\_cube\_le, symmExcess\_nonneg, symmExcess\_le\_absMean, symmExcess\_le\_one, one\_le\_absMean\_mul\_thirdAbsMoment}{Upper0423/Analytic/Moments.lean}

The excess bounds of Lemma~\ref{lem:eta}, in the equivalent formulation through
the ratio $r=\EE|X-X'|^3/(2\beta)$, are due to [XL] (item (iii) of
Section~\ref{sec:attribution}); the proof above is an independent and shorter
rederivation of the two inequalities actually needed.

\begin{proposition}[Feasibility: $\beta+\eta\ge2$]
  \label{prop:feasible}
  For every admissible $\mu$ we have $\EE|X-X'|^3\ge4$, equivalently
  $\beta+\eta\ge2$, with equality for the Rademacher law.
\end{proposition}

\begin{proof}[Proof sketch]
  For $x>y$ one has $\tfrac16(x-y)^3=\int_y^x(x-t)(t-y)\dd t$, so
  \[
    |x-y|^3
    =6\int_\RR\bigl[(x-t)_+(t-y)_+ + (y-t)_+(t-x)_+\bigr]\dd t ,
  \]
  and taking expectations,
  $\EE|X-X'|^3=12\int_\RR P(t)N(t)\dd t$ with $P(t)=\EE(X-t)_+$ and
  $N(t)=\EE(t-X)_+=P(t)+t$. Subtracting the Rademacher stop-loss profile
  $a_0(t)=\tfrac12(1-t)_+$ turns the deficit into a Gram form: with
  \[
    G(x,y)=\int_0^\infty\bigl((x-t)_+-a_0(t)\bigr)\bigl((y-t)_+-a_0(t)\bigr)\dd t ,
  \]
  one has the pointwise minorant $|x-y|^3\ge \text{(explicit)}+12\,G(x,y)$, and
  $\iint G\dd\mu\dd\mu\ge0$ by Fubini because $G$ is a positive semidefinite
  kernel (a square of an integral of a single function). The explicit part
  integrates to exactly~$4$.
\end{proof}
\formal{radStopLoss, stopLossGram, stopLossGram\_integral\_nonneg, four\_le\_integral\_symm\_abs\_cube, two\_le\_thirdAbsMoment\_add\_symmExcess}{Upper0423/Analytic/StopLoss.lean}

\begin{remark}
  The stop-loss formulation of Proposition~\ref{prop:feasible} is due to [XL]
  (item (iii) of Section~\ref{sec:attribution}); the proof above is a
  self-contained rederivation. Proposition~\ref{prop:feasible} is what makes the
  coordinate $a=\beta+\eta$ useful: it is bounded below by $2$ and it is
  precisely the quantity that controls the modulus of $f$ through the convex
  profile (Section~\ref{sec:kappa}).
\end{remark}

Summarizing, the parameter region used by the certificates is
\begin{equation}
  \label{eq:region}
  \beta\ge1,\qquad 0\le\eta\le s\le1,\qquad s\beta\ge1,\qquad \beta+\eta\ge2 .
\end{equation}
Two coordinate systems on \eqref{eq:region} are used:
\begin{itemize}[leftmargin=1.6em]
\item for $\beta\le2$: $\delta=\beta-1\in[0,1]$ and $w\in[0,1]$ with
  $\eta=1-\delta w$ (so $w$ interpolates between $\eta=1$ and $\eta=1-\delta$;
  note $\eta\ge 2-\beta=1-\delta$ by \eqref{eq:region});
\item for $\beta\ge2$: the pair $(\beta,\eta)$ directly, with $\eta\in[0,1]$.
\end{itemize}

\section{The Prawitz smoothing inequality}
\label{sec:prawitz}

\subsection{The kernel}

\begin{definition}
  \label{def:kernel}
  For $s\in(0,1)$ put
  \[
    k(s)=(1-s)\frac{\cos\pi s}{\sin\pi s}+\frac1\pi,
    \qquad
    K_\sigma(s)=\sigma(1-s)+\ii\,k(s)\qquad(\sigma\in\{\pm1\}) .
  \]
  We abbreviate $K=K_1$.
\end{definition}
\formal{prawitzSineKernel, prawitzKernel}{BerryEsseen/PrawitzKernel.lean, BerryEsseen/PrawitzFourier.lean}

\noindent Since $|K_{-1}(s)|=|K_1(s)|$ we only ever need
\begin{equation}
  \label{eq:absK}
  |K(s)|=\sqrt{(1-s)^2+k(s)^2} .
\end{equation}
Two equivalent rewritings of \eqref{eq:absK}, both used by the engine, follow
from the cotangent defect $q(z)=\tfrac1z-\cot z$:
\begin{align}
  \label{eq:K-right}
  |K(s)| &= (1-s)\sqrt{1+q(\pi(1-s))^2},\\
  \label{eq:K-left}
  s\,|K(s)| &= \sqrt{\bigl(s(1-s)\bigr)^2+\Bigl(\tfrac1\pi-s(1-s)\,q(\pi s)\Bigr)^2} .
\end{align}
Indeed $k(s)=(1-s)q(\pi(1-s))$ using $\cot(\pi-z)=-\cot z$, which gives
\eqref{eq:K-right}; and $k(s)=\tfrac1{\pi s}-(1-s)q(\pi s)$, which gives
\eqref{eq:K-left} after multiplying by $s$. Form \eqref{eq:K-left} is
numerically stable near $s=0$ (where $|K|$ blows up like $1/(\pi s)$ but
$s|K(s)|$ is smooth) and form \eqref{eq:K-right} near $s=1$. This is why the
high-frequency panels of a certificate are never allowed to straddle $s=1/2$: on
$[\tau,1/2]$ the engine uses \eqref{eq:K-left}, on $[1/2,1]$ it uses
\eqref{eq:K-right}.

Finally, the normal correction uses
\begin{equation}
  \label{eq:K-normal}
  \Bigl|K(s)-\frac{\ii}{\pi s}\Bigr| = (1-s)\sqrt{1+q(\pi s)^2},
\end{equation}
because $\tfrac1\pi-\tfrac1{\pi s}=-\tfrac{1-s}{\pi s}$, so the imaginary part of
$K(s)-\ii/(\pi s)$ equals $-(1-s)q(\pi s)$.
\formal{kernel\_norm\_eq\_sqrt, kernelR\_left, kernelR\_right, kernelR\_scaled, kernelR\_normal}{Upper0423/LargeN/Kernels.lean, Upper0423/Analytic/Kernels.lean}

\subsection{The smoothing bound}

\begin{definition}
  \label{def:smoothing}
  For a probability measure $\rho$ on $\RR$ with characteristic function
  $\varphi$, and for $T>0$, $\tau\in(0,1]$, put
  \begin{align*}
    \mathcal S(\rho,T,\tau)
    &= \int_0^\tau |K(t)|\;\bigl|\varphi(Tt)-e^{-T^2t^2/2}\bigr|\dd t
    && \text{(low)}\\
    &\quad+\int_\tau^1 |K(t)|\;\bigl|\varphi(Tt)\bigr|\dd t
    && \text{(high)}\\
    &\quad+\int_0^\tau \Bigl|K(t)-\frac{\ii}{\pi t}\Bigr| e^{-T^2t^2/2}\dd t
    && \text{(normal correction)}\\
    &\quad+\int_\tau^\infty \frac{e^{-T^2t^2/2}}{\pi t}\dd t .
    && \text{(tail)}
  \end{align*}
\end{definition}
\formal{smoothingIntegral, prawitzLowError, prawitzHighMagnitude, prawitzNormalCorrection, prawitzNormalTail}{BerryEsseen/PrawitzSmoothing.lean, BerryEsseen/SmoothingIntegrability.lean}

\begin{theorem}[Prawitz smoothing]
  \label{thm:prawitz}
  Let $\rho$ be a probability measure on $\RR$ with $\int|y|\dd\rho<\infty$, let
  $R$ be its distribution function, $T>0$ and $\tau\in(0,1]$. Then for every
  $x\in\RR$,
  \[
    \bigl|R(x)-\Phi(x)\bigr|\;\le\;\mathcal S(\rho,T,\tau) .
  \]
\end{theorem}
\formal{prawitz\_smoothing, abs\_prawitzFourier\_residual\_le, normalizedSum\_prawitz\_smoothing}{BerryEsseen/PrawitzSmoothing.lean, BerryEsseen/SmoothingMajorants.lean}

Theorem~\ref{thm:prawitz} is the inequality of Prawitz \cite{Prawitz1972}, in the
normalization of [XL] (item (ii) of Section~\ref{sec:attribution}). The
formalization proceeds through the Fourier functional
\[
  \mathcal P(\rho,\sigma,T,x)=\int_0^1
  \operatorname{Re}\Bigl[K_\sigma(t)\,e^{-\ii Ttx}\,\varphi(Tt)\Bigr]\dd t ,
\]
which satisfies $R(x)-\tfrac12\le\mathcal P(\rho,1,T,x)$ and
$\mathcal P(\rho,-1,T,x)\le R(x)-\tfrac12$ by Prawitz's two-sided inversion
formula (formalized through a Fubini argument and an explicit majorant/minorant
of the indicator). Comparing both sides with the Gaussian value
$\Phi(x)-\tfrac12$, replacing $\varphi$ by $e^{-T^2t^2/2}$ in the range
$t\le\tau$ and by $0$ in the range $t>\tau$, and using
$|K_{\pm1}|=|K|$, produces the four terms of Definition~\ref{def:smoothing}. The
third and fourth terms are the price of the substitution: the exact Gaussian
Fourier functional is recovered up to
$|K(t)-\ii/(\pi t)|$ on $[0,\tau]$ and the pure Gaussian tail beyond $\tau$,
the latter because $\int_0^\infty \frac{e^{-T^2t^2/2}}{\pi t}\sin(Ttx)\dd t$ is
the exact inversion integral for $\Phi$.

\subsection{The i.i.d.\ scaling}

For an admissible $\mu$ and $n\ge1$, write $\rho_n$ for the law of $S_n$. Then
$\varphi_n(v)=f(v/\sqrt n)^n$, so with the choice
\begin{equation}
  \label{eq:Tscale}
  T=c\sqrt n \qquad (c>0)
\end{equation}
we get $\varphi_n(Tt)=f(ct)^n$ and $e^{-T^2t^2/2}=g(ct)^n$ with
$g(u)=e^{-u^2/2}$. Thus every integrand in Definition~\ref{def:smoothing} becomes
a function of the \emph{per-summand} frequency
\begin{equation}
  \label{eq:u}
  u=c\,t ,
\end{equation}
raised to the power $n$, times a kernel factor depending on $t$ alone:
\begin{equation}
  \label{eq:integrands}
  \begin{aligned}
    \text{low}(t) &= |K(t)|\cdot\bigl|f(u)^n-g(u)^n\bigr|,\\
    \text{high}(t) &= |K(t)|\cdot|f(u)|^n,\\
    \text{normal}(t) &= \Bigl|K(t)-\frac{\ii}{\pi t}\Bigr|\cdot g(u)^n .
  \end{aligned}
\end{equation}
\formal{charFun\_normalizedSum\_scaled, prawitzLowError\_scaled, prawitzHighMagnitude\_scaled}{Upper0423/Small/Leaf.lean}

The tail term has a closed reduction: substituting $v=T^2t^2/2$,
\begin{equation}
  \label{eq:tail}
  \int_\tau^\infty\frac{e^{-T^2t^2/2}}{\pi t}\dd t
  = \frac1{2\pi}\int_{nc^2\tau^2/2}^\infty\frac{e^{-v}}{v}\dd v
  = \frac{E_1(nc^2\tau^2/2)}{2\pi} .
\end{equation}
\formal{prawitzNormalTail\_integral\_eq}{Upper0423/Small/Tail.lean}

\section{Characteristic-function estimates}
\label{sec:chf}

This section collects the pointwise estimates for $f$ that the certificates
consume. In Lean they are bundled into a single structure
\begin{leancode}
structure ChfFacts (f : ℝ → ℂ) (β η s : ℝ) : Prop
\end{leancode}
with eighteen fields, and the theorem \Lean{chfFacts} proves that every admissible
law satisfies it with the parameters \eqref{eq:coords}. The engine's soundness is
proved from \Lean{ChfFacts} alone, so this list is exactly the analytic interface
between the mathematics and the computation.
\formal{ChfFacts}{Upper0423/Analytic/ModulusBridge.lean}
\formal{chfFacts}{Upper0423/Analytic/ChfFactsProof.lean}

\subsection{The cubic cosine constant and the convex profile}
\label{sec:kappa}

\begin{lemma}
  \label{lem:kappa}
  For all $v\in\RR$, $\cos v-1+\tfrac{v^2}2\le \kap|v|^3$ with
  $\kap=49581/500000$.
\end{lemma}

\begin{proof}[Proof sketch]
  By symmetry assume $v\ge0$. For $v\le2$ the Taylor bound
  $\cos v\le1-\tfrac{v^2}2+\tfrac{v^4}{24}$ suffices since $v^4/24\le\kap v^3$
  for $v\le24\kap$, and $24\kap>2$. For $v\ge6$ the left side is at most
  $\tfrac{v^2}2$ and $\tfrac12\le\kap v$. The interval $[2,6]$ is covered by
  seven polynomial certificates, each reducing to the nonnegativity of an
  explicit polynomial on a dyadic subinterval after substituting the degree-$20$
  Taylor majorant of $\cos$.
\end{proof}
\formal{cos\_le\_taylor\_twenty, kappa\_gap\_nonneg\_0 \dots kappa\_gap\_nonneg\_6, cos\_sub\_one\_add\_half\_sq\_le\_kappa}{Upper0423/Analytic/CosineProfile.lean}

The constant $\kap$ and this use of it are due to [XL] (item (vii) of
Section~\ref{sec:attribution}); the true maximum is
$0.09916191350\ldots$, and $49581/500000=0.099162$ is a slightly coarser rational
majorant which is what all certificates are calibrated against.

Averaging Lemma~\ref{lem:kappa} over $v=uX$ gives at once:

\begin{proposition}
  \label{prop:rekappa}
  For all $u\in\RR$,
  $\operatorname{Re}f(u)\le1-\tfrac{u^2}2+\kap\beta|u|^3$, and
  $|f(u)|^2\le1-u^2+2\kap(\beta+\eta)|u|^3$.
\end{proposition}

\begin{proof}
  The first bound is Lemma~\ref{lem:kappa} with $v=uX$, using $\EE X^2=1$. For
  the second, note $|f(u)|^2=\EE\cos\bigl(u(X-X')\bigr)$ for an independent copy
  $X'$; apply Lemma~\ref{lem:kappa} with $v=u(X-X')$ and use
  $\EE(X-X')^2=2$ and $\EE|X-X'|^3=2(\beta+\eta)$.
\end{proof}
\formal{charFun\_re\_le\_kappa, normSq\_charFun\_le\_kappa}{Upper0423/Analytic/CosineProfile.lean}

The second bound of Proposition~\ref{prop:rekappa} degenerates for large $|u|$.
The fix, again due to [XL], is to replace the supporting line
$\tfrac12-\kap v$ by a globally convex profile.

\begin{definition}[Convex profile]
  \label{def:profile}
  With $A=(8-64\kap)/(1-\cos4)$ put
  \[
    q(v)=
    \begin{cases}
      \tfrac12-\kap v, & 0\le v\le 4,\\[2pt]
      A\,\dfrac{1-\cos v}{v^2}, & 4< v\le 2\pi,\\[6pt]
      0, & v>2\pi,
    \end{cases}
    \qquad q(v)=q(|v|) \text{ for } v<0 .
  \]
\end{definition}

\begin{lemma}
  \label{lem:profile}
  $q$ is continuous, convex on $\RR$, nonincreasing and nonnegative on
  $[0,\infty)$, and satisfies $\cos v\le 1-v^2q(|v|)$ for all $v$.
\end{lemma}

\begin{proof}[Proof sketch]
  The value $A$ is chosen so that the two branches agree at $v=4$: indeed
  $A(1-\cos4)=8-64\kap$ gives $A(1-\cos4)/16=\tfrac12-4\kap$. One checks
  $\tfrac{\dd}{\dd v}\bigl(A(1-\cos v)/v^2\bigr)\le0$ and
  $\tfrac{\dd^2}{\dd v^2}\bigl(A(1-\cos v)/v^2\bigr)\ge0$ on $[4,2\pi]$ by
  explicit sign estimates for $\sin$ and $\cos$ there, and that the right
  derivative at $4$ dominates $-\kap$; convexity of the splice follows, and the
  third branch is joined because $A(1-\cos2\pi)/(2\pi)^2=0$ with vanishing
  derivative. The inequality $\cos v\le1-v^2q(|v|)$ holds on $[0,4]$ by
  Lemma~\ref{lem:kappa}, is an equality on $(4,2\pi]$ by definition of $A$
  (up to the factor $A\le1$), and is trivial beyond $2\pi$.
\end{proof}
\formal{profile, convexOn\_profile, antitone\_profile, cos\_le\_profile}{Upper0423/Analytic/CosineProfile.lean}

\begin{proposition}[Global modulus bound]
  \label{prop:profilebound}
  For all $u$,
  \[
    |f(u)|^2\le1-2u^2\,q\bigl((\beta+\eta)|u|\bigr) .
  \]
  In particular, if $\beta+\eta\le D$ and $4\le D|u|\le2\pi$, then
  \[
    |f(u)|^2\le1-\frac{2A\bigl(1-\cos(D|u|)\bigr)}{D^2} .
  \]
\end{proposition}

\begin{proof}[Proof sketch]
  Write $|f(u)|^2=\EE\cos\bigl(u(X-X')\bigr)\le 1-u^2\,\EE\bigl[(X-X')^2
  q(|u(X-X')|)\bigr]$ by Lemma~\ref{lem:profile}. The measure
  $\tfrac12(X-X')^2\dd\PP$ has total mass $1$ and, by convexity of $q$ and
  Jensen's inequality against that weighted measure,
  $\EE\bigl[\tfrac12(X-X')^2q(|u||X-X'|)\bigr]\ge
  q\bigl(|u|\,\EE[\tfrac12(X-X')^2|X-X'|]\bigr)
  = q\bigl(|u|(\beta+\eta)\bigr)$, using
  $\EE|X-X'|^3=2(\beta+\eta)$. The second display uses that $q$ is
  nonincreasing, so $q(D|u|)\le q((\beta+\eta)|u|)$ is false in the wrong
  direction --- one applies monotonicity in the correct direction, replacing
  $(\beta+\eta)|u|$ by the larger $D|u|$ and using the antitonicity of $q$ on
  $[0,\infty)$ together with the explicit curved branch.
\end{proof}
\formal{normSq\_charFun\_le\_profile, normSq\_charFun\_le\_curved}{Upper0423/Analytic/CosineProfile.lean, Upper0423/Analytic/ChfFactsProof.lean}

\subsection{The real part}

\begin{proposition}[Signed real lower bound]
  \label{prop:relower}
  For $0\le u\le\pi$, $\cos u\le\operatorname{Re}f(u)$.
\end{proposition}

\begin{proof}
  Fix $u\in(0,\pi]$ and put $c=\sin u/(2u)$. We claim
  $H(v)=\cos v+cv^2$ satisfies $H(v)\ge H(u)$ for all $v\ge0$. The chord bound
  $w\sin u\le u\sin w$ for $0\le w\le u\le\pi$ (concavity of $\sin$ on $[0,\pi]$)
  gives $H'(v)=2cv-\sin v\le 0$ for $v\le u$ and $\ge0$ for $v\ge u$ in the
  relevant range; the case $v>\pi$ is handled directly since $\sin v\le 1$ and
  $2cv\ge 2cu=\sin u$. Substituting $v=u|X|$ and taking expectations, using
  $\EE X^2=1$ and $\operatorname{Re}f(u)=\EE\cos(uX)=\EE\cos(u|X|)$, yields
  $\operatorname{Re}f(u)+c\ge\cos u+cu^2\cdot 1$ --- i.e.
  $\operatorname{Re}f(u)\ge\cos u$.
\end{proof}
\formal{mul\_sin\_le\_mul\_sin, cos\_add\_quad\_le, cos\_le\_charFun\_re}{Upper0423/Analytic/RealLower.lean}

No bounded-support hypothesis is needed. This estimate is what allows the engine
to bracket $\operatorname{Re}f$ from both sides and hence to use the
\emph{signed} real part, rather than only its modulus, in the coordinate modulus
branch (Section~\ref{sec:modulus}).

\begin{proposition}[Real upper bounds]
  \label{prop:reupper}
  For $u\ge0$ we have the three bounds
  \begin{align}
    \operatorname{Re}f(u)&\le \cos u+\frac{(\beta-1)u^3}{6},
      \label{eq:re-anchor}\\
    \operatorname{Re}f(u)&\le \cos u+(1-s)\bigl(u\sin u+u^2\bigr),
      \label{eq:re-radial}\\
    \operatorname{Re}f(u)&\le 1-\frac{u^2}{2}+\kap\beta u^3 .
      \label{eq:re-kappa}
  \end{align}
\end{proposition}

Bound \eqref{eq:re-anchor} is the two-sided anchor
$|\operatorname{Re}f(u)-\cos u|\le(\beta-1)|u|^3/6$, which follows from
$|\cos(uy)-\cos u-\tfrac12(1-y^2)(\ldots)|$-type Taylor estimates weighted by
$\EE\bigl||X|^3-\ldots\bigr|$; \eqref{eq:re-radial} is the radial bound of
Proposition~\ref{prop:radial}; \eqref{eq:re-kappa} is
Proposition~\ref{prop:rekappa}.

\subsection{The imaginary part}

Four bounds are used; all four are needed, because they are sharp in different
corners of the region \eqref{eq:region}.

\begin{proposition}[Anchored cubic]
  \label{prop:imcubic}
  For all $u$,
  \[
    |\operatorname{Im}f(u)|\le \frac{\sqrt{(\beta-1)(\beta+5/3)}}{6}\,|u|^3 .
  \]
\end{proposition}

\begin{proof}[Proof sketch]
  $\operatorname{Im}f(u)=\EE\sin(uX)$ and $\EE(uX)=0$, so
  $\operatorname{Im}f(u)=\EE[\sin(uX)-uX]$ and
  $|\sin v-v|\le|v|^3/6$ gives the crude $\beta|u|^3/6$. The improvement comes
  from noticing that the relevant moment is
  $\EE\bigl||X|\,(X^2-1)\bigr|$ rather than $\EE|X|^3$: by Cauchy--Schwarz with
  the weight $w(y)=(y-1)^2(2y+1)$ of Lemma~\ref{lem:betapos},
  \[
    \bigl(\EE|X|\,|X^2-1|\bigr)^2\le \EE w(|X|)\cdot\EE\frac{\ldots}{\ldots}
    \le (\beta-1)(\beta+5/3),
  \]
  the second inequality being an explicit polynomial comparison. The anchoring
  at $\cos u$ (rather than at $1$) is what makes the moment $(\beta-1)$-small.
\end{proof}
\formal{charFun\_im\_cubic\_anchor}{BerryEsseen/ImaginaryPart.lean (via Upper0423/Analytic/ChfFactsProof.lean)}

\begin{proposition}[Sharpened quadratic]
  \label{prop:imquad}
  For all $u$,
  $|\operatorname{Im}f(u)|\le\dfrac{\sqrt{2(\beta-1)/3}}{2}\,u^2$.
\end{proposition}

\begin{proof}
  The scalar inequality $|\sin(uy)-y\sin u|\le\tfrac{u^2}{2}\,y|y-1|$ for
  $y\ge0$, applied with $y=|X|$ and combined with $\EE X=0$, gives
  $|\operatorname{Im}f(u)|\le\tfrac{u^2}{2}\,\EE\bigl[|X|\,\bigl||X|-1\bigr|\bigr]$.
  The polynomial identity
  \[
    w(y)\frac{y^2+2y}{9}-y^2(y-1)^2 = \frac{2y(y-1)^4}{9}\ \ge 0
    \qquad (y\ge0),
  \]
  with $w(y)=(y-1)^2(2y+1)$ and $\EE w(|X|)=2(\beta-1)$, gives by weighted
  Cauchy--Schwarz
  \[
    \Bigl(\EE\bigl[|X|\,\bigl||X|-1\bigr|\bigr]\Bigr)^2
    \le 2(\beta-1)\,\frac{1+2s}{9}\le\frac{2(\beta-1)}{3},
  \]
  using $s\le1$.
\end{proof}
\formal{weighted\_linear\_sq\_le\_sharp, integral\_abs\_weighted\_linear\_bound\_sharp, abs\_charFun\_im\_le\_quadratic\_sharp}{Upper0423/Analytic/QuadraticImag.lean}

\begin{proposition}[Sine circle]
  \label{prop:circle}
  For all $u$,
  \[
    |\operatorname{Im}f(u)|\le\frac{\sqrt{\beta^2-\eta^2}}{6}\,|u|^3 .
  \]
\end{proposition}

\begin{proof}[Proof sketch]
  This is the two-point decomposition of [XL] (item (iv) of
  Section~\ref{sec:attribution}). Write $r_0=\EE X_+=s/2$ and decompose $\mu$ as
  a mixture of centered two-point laws $P_{a,b}$ (mass $b/(a+b)$ at $a$, mass
  $a/(a+b)$ at $-b$) with mixing measure
  $\dd\nu=\tfrac{a+b}{r_0}\dd\mu_+(a)\dd\mu_-(b)$. For each component, with
  $h_{a,b}(u)=\operatorname{Im}\widehat{P_{a,b}}(u)$,
  \[
    |h_{a,b}(u)|\le\frac{|u|^3\,ab\,|a^2-b^2|}{6(a+b)},
  \]
  and with $\beta_{a,b}=\frac{ab(a^2+b^2)}{a+b}$ and
  $e_{a,b}=\frac{2a^2b^2}{a+b}$ one has the pointwise \emph{circle} inequality
  \[
    e_{a,b}^2+\Bigl(\frac{6\,h_{a,b}(u)}{|u|^3}\Bigr)^2\le \beta_{a,b}^2 .
  \]
  Minkowski's inequality in $L^2(\nu)$ upgrades this to
  $\zeta^2+\bigl(6\operatorname{Im}f(u)/|u|^3\bigr)^2\le\beta^2$ where
  $\zeta=\int e_{a,b}\dd\nu$, and an explicit computation gives
  $\zeta=4A_2B_2/s\ge s\ge\eta$ (with $A_2=\EE X_+^2$, $B_2=\EE X_-^2$), using
  Lemma~\ref{lem:eta}. Hence
  $\bigl(6\operatorname{Im}f(u)/|u|^3\bigr)^2\le\beta^2-\eta^2$.
\end{proof}
\formal{abs\_sin\_twoPoint\_le, twoPoint\_circle, sineCircle\_prod, abs\_charFun\_im\_le\_sineCircle}{Upper0423/Analytic/SineCircle.lean}

\begin{proposition}[Radial disk]
  \label{prop:radial}
  Put $d=1-s$ and $Y=|X|$, $A=\operatorname{sgn}X$ (so $X=AY$), $a=\EE A$. For
  $0\le u\le\pi$,
  \begin{align*}
    \bigl|f(u)-\bigl[\cos u+d\,u\sin u+\ii\,a(\sin u-u\cos u)\bigr]\bigr|
      &\le d\,u^2,\\
    \operatorname{Re}f(u)&\le\cos u+d\,(u\sin u+u^2),\\
    |\operatorname{Im}f(u)|&\le\sqrt{2d}\,(\sin u-u\cos u)+d\,u^2 .
  \end{align*}
\end{proposition}

\begin{proof}[Proof sketch]
  The second-order remainder bound
  $\bigl|e^{\ii uY}-e^{\ii u}\bigl(1+\ii u(Y-1)\bigr)\bigr|\le\tfrac{u^2}2(Y-1)^2$
  and $\EE(Y-1)^2=2-2s=2d$ give the first display after averaging, with the
  $\ii a$ coefficient arising from $\EE[A\,(\ldots)]$ and $\EE[AY]=\EE X=0$.
  Cauchy--Schwarz applied to $0=\EE[AY]$ gives $a^2+s^2\le1$, hence
  $a^2\le1-s^2\le2d$, which yields the third display; the second follows from
  the real part of the first together with $\sin u-u\cos u\ge0$ on $[0,\pi]$.
\end{proof}
\formal{abs\_integral\_radialSign\_le, charFun\_re\_le\_radial, abs\_charFun\_im\_le\_radial}{Upper0423/Analytic/Radial.lean}

Proposition~\ref{prop:radial} is the estimate of [XL] (item (v) of
Section~\ref{sec:attribution}); the form above retains the complex remainder norm
and the sharpened sign-mean estimate $a^2\le 1-s^2$.

\subsection{The joint anchor disk}

\begin{proposition}[ODE anchor]
  \label{prop:ode}
  For $0\le u\le\pi$,
  \[
    \bigl\|f(u)-\cos u\bigr\|\le \sqrt{(\beta-1)(\beta+5/3)}\;(u-\sin u) .
  \]
\end{proposition}

\begin{proof}[Proof sketch]
  For fixed real $x$ the function
  \[
    G(u)=e^{\ii ux}-\cos u-\ii x\sin u-(1-x^2)(1-\cos u)
  \]
  satisfies $G(0)=G'(0)=0$ and $G''+G=(1-x^2)(e^{\ii ux}-1)$, whence
  $G(u)=\int_0^u\sin(u-v)\,(1-x^2)(e^{\ii vx}-1)\dd v$. Using
  $|e^{\ii vx}-1|\le v|x|$ and $\int_0^u v\sin(u-v)\dd v=u-\sin u$,
  \[
    |G(u)|\le |x|\,|x^2-1|\,(u-\sin u) .
  \]
  Averaging over $x=X$ kills the terms $\ii x\sin u$ and $(1-x^2)(1-\cos u)$
  because $\EE X=0$ and $\EE X^2=1$, leaving
  $\|f(u)-\cos u\|\le \EE\bigl[|X|\,|X^2-1|\bigr](u-\sin u)$, and the moment is
  bounded by $\sqrt{(\beta-1)(\beta+5/3)}$ as in
  Proposition~\ref{prop:imcubic}.
\end{proof}
\formal{odeAnchorRemainder, norm\_odeAnchorRemainder\_le, norm\_charFun\_sub\_cos\_le\_ode}{Upper0423/Analytic/OdeAnchor.lean}

Finally the plain Taylor disk is retained:
\begin{equation}
  \label{eq:quadratic}
  \Bigl\|f(u)-1+\frac{u^2}{2}\Bigr\|\le\frac{\beta|u|^3}{6} .
\end{equation}

\subsection{Summary of the interface}

Table~\ref{tab:chf} lists the eighteen fields of \Lean{ChfFacts}. Each is proved
for admissible laws in the file indicated.

\begin{table}[htbp]
  \centering\footnotesize
  \begin{tabular}{llp{5.2cm}}
    \toprule
    Field & Content & Source file\\
    \midrule
    \Lean{one\_le\_beta} & $1\le\beta$ & \file{MomentExcess.lean}\\
    \Lean{eta\_nonneg}, \Lean{eta\_le\_s}, \Lean{s\_le\_one}, \Lean{s\_beta}
      & $0\le\eta\le s\le1$, $s\beta\ge1$ & \file{Analytic/Moments.lean}\\
    \Lean{norm\_le\_one} & $|f|\le1$ & Mathlib\\
    \Lean{re\_ge\_cos} & Prop.~\ref{prop:relower} & \file{Analytic/RealLower.lean}\\
    \Lean{re\_anchor} & \eqref{eq:re-anchor} & \file{RealPart.lean}\\
    \Lean{re\_radial} & \eqref{eq:re-radial} & \file{Analytic/Radial.lean}\\
    \Lean{re\_kappa} & \eqref{eq:re-kappa} & \file{Analytic/CosineProfile.lean}\\
    \Lean{im\_cubic} & Prop.~\ref{prop:imcubic} & \file{ImaginaryPart.lean}\\
    \Lean{im\_quad} & Prop.~\ref{prop:imquad} & \file{Analytic/QuadraticImag.lean}\\
    \Lean{im\_circle} & Prop.~\ref{prop:circle} & \file{Analytic/SineCircle.lean}\\
    \Lean{im\_radial} & Prop.~\ref{prop:radial} & \file{Analytic/Radial.lean}\\
    \Lean{normSq\_kappa} & Prop.~\ref{prop:rekappa} & \file{Analytic/CosineProfile.lean}\\
    \Lean{normSq\_curved} & Prop.~\ref{prop:profilebound} & \file{Analytic/ChfFactsProof.lean}\\
    \Lean{ode} & Prop.~\ref{prop:ode} & \file{Analytic/OdeAnchor.lean}\\
    \Lean{quadratic} & \eqref{eq:quadratic} & \file{CharacteristicBounds.lean}\\
    \bottomrule
  \end{tabular}
  \caption{The analytic interface \Lean{ChfFacts}.}
  \label{tab:chf}
\end{table}

\section{The case decomposition}
\label{sec:split}

The proof of Theorem~\ref{thm:main} splits the parameter space
$\{(\mu,n,x)\}$ into four regions, as follows
(\Lean{normalizedError\_le\_0423} in \file{Upper0423/Main.lean}):

\begin{center}
\footnotesize
\setlength{\tabcolsep}{4pt}
\begin{tabular}{llll}
  \toprule
  Region & Lean theorem & File & Method\\
  \midrule
  $n=1$ & \Lean{normalizedError\_one\_le} & \file{SingleSummand.lean} & polynomial majorants\\
  $\ell\ge\frac{13003}{10000}$ & \Lean{normalizedError\_large\_ell} & \file{LargeEll.lean} & Cantelli\\
  $2\le n\le255$, $\ell<\frac{13003}{10000}$ & \Lean{normalizedError\_moderate\_le} & \file{Covers/All.lean} & $569$ covers\\
  $n\ge256$, $\ell\le\frac{13003}{10000}$ & \Lean{normalizedError\_uniform} & \file{Uniform/Data/All.lean} & $281$ cells\\
  \bottomrule
\end{tabular}
\end{center}

\noindent Here $\ell=\beta/\sqrt n$. The four cases are exhaustive:
if $n=1$ the first applies; otherwise either $\ell\ge13003/10000$ (second case)
or $\ell<13003/10000$, and then the third or fourth applies according to
$n\le255$ or $n\ge256$. The overlap of the second and fourth cases at
$\ell=13003/10000$ is harmless (both are stated with the appropriate non-strict
or strict inequality).

The threshold $13003/10000$ is exactly what the Cantelli cap allows; see
Section~\ref{sec:largeell}. Its role in the finite covers is that for
$2\le n\le255$ one needs only
\[
  1\le\beta\le \frac{13003}{10000}\sqrt n ,\qquad 0\le\eta\le1 ,
\]
a compact rational box for each $n$.

\section{The single summand: $n=1$}
\label{sec:one}

\begin{theorem}
  \label{thm:one}
  For every admissible $\mu$ and every $t\in\RR$,
  \[
    \bigl|\PP(X\le t)-\Phi(t)\bigr|\le\frac{21}{50}\,\beta .
  \]
  Consequently $\mathcal E(\mu,1,x)\le 21/50\le 423/1000$.
\end{theorem}

Since $21/50=0.42<0.423$, this case needs no computation at all. The mechanism is
a family of \emph{rational polynomial majorants} of indicator functions.

\begin{lemma}[Majorant certificate]
  \label{lem:maj}
  Let $C,a,b,c,r\in\QQ$ and suppose
  \[
    \varphi(y):=C|y|^3+ay+by^2+c
  \]
  satisfies $\varphi\ge0$ on $\RR$ and $\varphi\ge1$ on $(-\infty,r]$. Then for
  every admissible $\mu$,
  \[
    \PP(X\le r)\le \EE\varphi(X)= C\beta+b+c ,
  \]
  using $\EE X=0$ and $\EE X^2=1$.
\end{lemma}

\begin{proof}
  $\mathbf 1_{(-\infty,r]}\le\varphi$ pointwise, and $\varphi$ is $\mu$-integrable
  by admissibility; the displayed value of $\EE\varphi$ is immediate.
\end{proof}
\formal{prob\_le\_le\_integral, integral\_cubic\_majorant, prob\_le\_le\_of\_certificate}{Upper0423/SingleSummand.lean}

The hypotheses of Lemma~\ref{lem:maj} are themselves certified by two
sum-of-squares identities, which is what makes them checkable by \Lean{nlinarith}
on rational data.

\begin{lemma}[SOS1]
  \label{lem:sos1}
  Let $C>0$ and $v\in\RR$ with $b+2Cv>0$ and
  $(Cv^2-a)^2\le4(b+2Cv)c$. Then $Cz^3+bz^2+az+c\ge0$ for all $z\ge0$.
\end{lemma}
\begin{proof}
  Multiply by $4(b+2Cv)>0$ and use the identity
  \begin{multline*}
    4(b+2Cv)\bigl(Cz^3+bz^2+az+c\bigr)
    =4C(b+2Cv)\,z(z-v)^2\\
    +\bigl(2(b+2Cv)z-(Cv^2-a)\bigr)^2
    +\bigl(4(b+2Cv)c-(Cv^2-a)^2\bigr) .
  \end{multline*}
  Each of the three summands is nonnegative.
\end{proof}

\begin{lemma}[SOS2]
  \label{lem:sos2}
  Let $C>0$ and $b^2\le-3Ca$. Then $z\mapsto Cz^3+bz^2-az+c$ is nondecreasing.
\end{lemma}
\begin{proof}
  For $w\le z$,
  \begin{multline*}
    12C\Bigl[\bigl(Cz^3+bz^2-az+c\bigr)-\bigl(Cw^3+bw^2-aw+c\bigr)\Bigr]\\
    =(z-w)\Bigl(3C^2(z-w)^2+\bigl(3C(z+w)+2b\bigr)^2+4\bigl(-3Ca-b^2\bigr)\Bigr)\ge0 .
  \end{multline*}
\end{proof}
\formal{cubic\_nonneg, cubic\_mono, certificate\_nonpos, certificate\_last}{Upper0423/SingleSummand.lean}

\begin{proof}[Proof of Theorem~\ref{thm:one}]
  By replacing $X$ with $-X$ it suffices to prove the one-sided bound
  $\PP(X\le t)-\Phi(t)\le\tfrac{21}{50}\beta$ for all $t$ (the reverse
  discrepancy is obtained from the one-sided bound for $-X$ together with
  $\Phi(-t)=1-\Phi(t)$, and $\beta$ is unchanged). Recall $\beta\ge1$
  (Lemma~\ref{lem:betapos}). The range of $t$ is split into ten pieces:
  \begin{itemize}[leftmargin=1.6em]
  \item $t\le-6/5$: a pure quadratic majorant ($C=0$) with
    $(a,b,c)=\bigl(-\tfrac53\cdot\tfrac{900}{3721},\tfrac{900}{3721},
    \tfrac{25}{36}\cdot\tfrac{900}{3721}\bigr)$ and $r=-6/5$ gives
    $\PP(X\le t)\le b+c$, which is below $\tfrac{21}{50}\beta$ already since
    $\Phi\ge0$.
  \item $-6/5<t\le-1$: another pure quadratic majorant with
    $(a,b,c)=(-\tfrac12,\tfrac14,\tfrac14)$, $r=-1$, combined with the explicit
    Gaussian lower bound of Lemma~\ref{lem:gauss} at $z=6/5$.
  \item seven cells with endpoints
    $-1,\,-\tfrac9{10},\,-\tfrac45,\,-\tfrac35,\,-\tfrac25,\,-\tfrac15,\,0,\,\tfrac3{10}$:
    each carries a cubic certificate with $C=\tfrac{21}{50}$ and explicit
    rational $(a,b,c,v)$, verified by Lemmas~\ref{lem:sos1} and \ref{lem:sos2},
    together with the Gaussian value at the left endpoint of the cell.
  \item $t>3/10$: $\PP(X\le t)\le1$ and
    $\Phi(t)\ge\Phi(3/10)\ge\tfrac12+\tfrac{197}{1680}$, which suffices.
  \end{itemize}
  In each of the seven cells the estimate has the shape
  \[
    \PP(X\le t)-\Phi(t)
    \le \underbrace{\PP(X\le r)}_{\le\,C\beta+b+c}-\underbrace{\Phi(l)}_{\ge\,b+c}
    \le C\beta ,
  \]
  where $[l,r]$ is the cell, using monotonicity of both $F$ and $\Phi$; the
  numerical content of each row is the pair of inequalities $b+c\le\Phi(l)$ and
  the SOS conditions.
\end{proof}
\formal{prob\_le\_sub\_normalCDF\_le, abs\_prob\_le\_sub\_normalCDF\_le, normalizedError\_one\_le}{Upper0423/SingleSummand.lean}

\begin{lemma}[Explicit Gaussian bounds]
  \label{lem:gauss}
  For $z\ge0$,
  $\Phi(-z)\ge\tfrac12-\tfrac25\bigl(z-\tfrac{z^3}{6}+\tfrac{z^5}{40}\bigr)$, and
  $\Phi(3/10)\ge\tfrac12+\tfrac{197}{1680}$.
\end{lemma}
\begin{proof}
  From $e^{-v}\le1-v+\tfrac{v^2}2$ for $v\ge0$ one gets
  $\phi(x)\le\tfrac25-\tfrac{x^2}5+\tfrac{x^4}{20}$; integrating over $[-z,0]$
  gives the first bound. From $1-\tfrac{x^2}2\le e^{-x^2/2}$ and the numerical
  bound $(2\pi)^{-1/2}\ge25/63$ one gets
  $\phi(x)\ge\tfrac{25}{63}(1-\tfrac{x^2}2)$ on $[0,3/10]$; integrating gives the
  second.
\end{proof}
\formal{exp\_neg\_le\_quadratic, gaussianPDFReal\_le\_quartic, normalCDF\_neg\_ge, normalCDF\_three\_tenths\_ge}{Upper0423/SingleSummand.lean}

\section{The large Lyapunov fraction}
\label{sec:largeell}

\begin{theorem}
  \label{thm:cantelli}
  For every admissible $\mu$, every $n\ge1$ and every $x$,
  $\bigl|F_n(x)-\Phi(x)\bigr|\le\tfrac{11}{20}$. Consequently, if
  $\ell=\beta/\sqrt n\ge13003/10000$, then
  \[
    \mathcal E(\mu,n,x)=\frac{\sqrt n}{\beta}\bigl|F_n(x)-\Phi(x)\bigr|
    \le\frac{10000}{13003}\cdot\frac{11}{20}=\frac{5500}{13003}
    <\frac{423}{1000} .
  \]
\end{theorem}

\begin{proof}
  $S_n$ has mean $0$ and variance $1$, so Cantelli's (one-sided Chebyshev)
  inequality gives $\PP(S_n\ge x)\le\tfrac{1}{1+x^2}$ for $x>0$ and hence
  $F_n(x)\ge\tfrac{x^2}{1+x^2}$; symmetrically
  $F_n(x)\le\tfrac{1}{1+x^2}$ for $x<0$. For $x\ge0$ one then needs the
  elementary Gaussian gap estimate
  \[
    \Phi(x)-\frac{x^2}{1+x^2}\le\frac{11}{20}\qquad(x\ge0),
  \]
  proved in \file{GaussianBounds.lean} by a polynomial argument on $[0,1]$ and by
  monotonicity beyond, together with $F_n\le1$ and $\Phi\ge\tfrac12$; the case
  $x<0$ is symmetric.
\end{proof}
\formal{cantelli\_upper, cantelli\_lower, normalCDF\_cantelli\_gap, cdf\_error\_le\_eleven\_twentieths, normalizedSum\_cdf\_error\_le\_eleven\_twentieths}{BerryEsseen/Cantelli.lean, BerryEsseen/GaussianBounds.lean}
\formal{normalizedError\_large\_ell}{Upper0423/LargeEll.lean}

\begin{remark}
  The threshold is determined by the numerical identity
  $\tfrac{11}{20}\cdot\tfrac{10000}{13003}=0.4229793\ldots<0.423$: any smaller
  threshold would fail, and $13003/10000$ is the value the finite covers are
  built for. A sharper unconditional cap than $11/20$ would lower the threshold
  and shrink the finite work; conversely, one cannot reuse a threshold calibrated
  for a weaker target, e.g. $56/45=1.2\overline{4}$, because
  $\tfrac{11}{20}\cdot\tfrac{45}{56}=0.4419\ldots>0.423$.
\end{remark}

\section{The interval and Taylor-jet engine}
\label{sec:engine}

The remaining two regions require numerical evaluation of the four integrals of
Definition~\ref{def:smoothing}. This section describes the machinery: an
executable fixed-point interval arithmetic with fourth-order Taylor jets, a
quadrature rule, a majorant assembly, and the Lean theorems that make the whole
thing sound. Nothing in this section is trusted; everything is proved.

\subsection{Fixed-point intervals}

\begin{definition}
  \label{def:iv}
  Let $S=2^{40}=1099511627776$. An \emph{interval} is a pair of integers
  $[\,\ell,h\,]$, and it \emph{contains} $x\in\RR$ when
  $\ell\le Sx\le h$.
\end{definition}
\formal{Iv, Iv.Contains, scale, precision}{Upper0423/Native/Interval.lean, Upper0423/Native/IntervalSound.lean}

All arithmetic is exact integer arithmetic with outward rounding, so containment
is preserved. Writing $\lfloor\cdot\rfloor$ and $\lceil\cdot\rceil$ for
floor/ceiling division of integers (defined for divisors of either sign), the
operations are
\[
  \begin{array}{ll}
  \text{negation} & [-h,-\ell];\\
  \text{addition} & [\ell_1+\ell_2,\;h_1+h_2];\\
  \text{subtraction} & [\ell_1-h_2,\;h_1-\ell_2];\\
  \text{multiplication} & \bigl[\lfloor\min(\ldots)/S\rfloor,\;\lceil\max(\ldots)/S\rceil\bigr]
    \text{ over the four endpoint products};\\
  \text{division by } n\in\NN & [\lfloor\ell/n\rfloor,\lceil h/n\rceil];\\
  \text{reciprocal} & \bigl[\lfloor S^2/h\rfloor,\lceil S^2/\ell\rceil\bigr]
    \text{ (fails if } \ell\le0\le h);\\
  \text{square} & \text{as multiplication, with lower bound } 0 \text{ if } \ell\le0\le h;\\
  \text{square root} & \bigl[\,\lfloor\sqrt{\ell S}\,\rfloor,\;\lceil\sqrt{hS}\,\rceil\,\bigr]
    \text{ via } \Lean{Nat.sqrt} \text{ (fails if }\ell<0).
  \end{array}
\]
Each of these has a soundness lemma of the form ``if $a$ contains $x$ and $b$
contains $y$ then $a\ast b$ contains $x\ast y$''. Exponentiation to a natural
power is by the same binary-exponentiation loop as the reference C\texttt{++}
implementation, so that the Lean checker and the search program agree bit for
bit.

\subsection{Transcendental seeds}

\begin{lemma}[$\pi$]
  \label{lem:pi}
  The interval $[3454217652352,\,3454217652372]$ contains $\pi$.
\end{lemma}
\begin{proof}
  Immediate from Mathlib's \Lean{Real.pi\_gt\_d20} and \Lean{Real.pi\_lt\_d20}.
\end{proof}
\formal{piIv, contains\_piIv}{Upper0423/Native/TranscendentalSound.lean}

\begin{lemma}[exponential]
  \label{lem:exp}
  The routine \Lean{Iv.exp} returns an interval containing $e^x$ whenever it
  succeeds.
\end{lemma}
\begin{proof}[Proof sketch]
  Endpointwise. For $x\le-40S$ return $[0,1]$, which is valid because
  $e^{-40}\le1/S$. Otherwise find the least $k$ with $8|x|\le S2^k$, put
  $y=x/2^k$ (so $|y|\le S/8$), sum twelve Taylor terms of $e^y$ by the recursion
  $(\Sigma,T)\mapsto(\Sigma+Ty/(j{+}1),\,Ty/(j{+}1))$, widen the result by one
  unit in each direction to absorb the truncation error, clamp at $0$, and square
  $k$ times. The Taylor remainder after twelve terms at $|y|\le1/8$ is far below
  one unit in the last place.
\end{proof}

\begin{lemma}[sine and cosine]
  \label{lem:sincos}
  The routine \Lean{Iv.sinCos} returns a pair of intervals containing
  $(\sin x,\cos x)$ for every $x$ in the argument interval.
\end{lemma}
\begin{proof}[Proof sketch]
  A point evaluation shifts the argument by the nearest multiple of $\pi/2$ into
  $[-1,1]$ (using the enclosure of Lemma~\ref{lem:pi}), evaluates
  degree-$17$ Taylor polynomials, widens by one unit, and permutes/negates
  according to the quadrant. For an interval argument, the two endpoint
  enclosures are hulled, and additionally hulled with $\pm1$ at every stationary
  point ($k\pi$ for $\cos$, $(2k{+}1)\pi/2$ for $\sin$, $|k|\le4$) whose own
  enclosure meets the argument interval; the result is intersected with
  $[-1,1]$. Soundness rests on the observation that between two consecutive
  stationary points the function is monotone, so its range over a subinterval is
  spanned by the endpoint values together with the extreme values at any
  interior stationary point. Arguments of absolute value larger than $7$ return
  $[-1,1]$.
\end{proof}
\formal{expEndpoint\_sound, Contains.exp, sinCosEndpoint\_sound, Contains.sinCos}{Upper0423/Native/TranscendentalSound.lean, Upper0423/Native/TrigSound.lean}

\subsection{Fourth-order jets}

\begin{definition}
  \label{def:jet}
  A \emph{jet} is a quintuple of intervals $(c_0,\dots,c_4)$. It
  \emph{encloses} a function $F$ at a point $\sigma$ when $c_k$ contains
  $F^{(k)}(\sigma)/k!$ for $k=0,\dots,4$. It encloses $F$ \emph{on} a set $D$
  when $F$ is four times differentiable on $D$ and the enclosure holds at every
  point of $D$.
\end{definition}
\formal{Jet, Jet.Encloses, Jet.EnclosesOn}{Upper0423/Native/JetSound.lean}

Jets are propagated by the usual normalized-coefficient recurrences: products by
the Cauchy convolution, reciprocal and square root by the Newton-type
recurrences, exponential and sine/cosine by the weighted convolutions
$r_k=\tfrac1k\sum_{j=1}^{k}j\,a_j r_{k-j}$ (and its trigonometric pair). The
zeroth coefficient of a squaring uses the dedicated interval square (so that
$c_0\ge0$ is exploited). Each of these has a Lean soundness lemma, reducing to
scale-free coefficient identities in the \Lean{BerryEsseen.Numerics} namespace;
only interval containment is redone in the engine's own namespace.

\subsection{The quadrature rule}
\label{sec:panel}

A \emph{panel} is a rational interval $[l,r]$. Parametrize it by
$t=m+h\sigma$ with $m=\tfrac{l+r}2$, $h=\tfrac{r-l}2$ and $\sigma\in[-1,1]$.

\begin{lemma}[Panel rule]
  \label{lem:panel}
  Let $F$ be a real function, $l<r$ rationals. Suppose a jet $W$ encloses
  $\sigma\mapsto F(m+h\sigma)$ on $[-1,1]$ and a jet $M$ encloses it at
  $\sigma=0$. Then
  \[
    \int_l^r F(t)\dd t\;\le\;
    \min\Bigl\{
      (r-l)\Bigl(\overline{M_0}+\frac{\overline{M_2}}{3}+\frac{\overline{W_4}}{5}\Bigr),\;
      (r-l)\,\overline{W_0}
    \Bigr\},
  \]
  where $\overline{\,\cdot\,}$ denotes the upper endpoint of an interval,
  rescaled by $S^{-1}$.
\end{lemma}
\begin{proof}[Proof sketch]
  For the second (range) bound, $F\le \overline{W_0}$ on the panel. For the first
  (Taylor) bound, expand
  $F(m+h\sigma)=\sum_{k\le3}c_k\sigma^k+c_4(\xi)\sigma^4$ with
  $c_k=c_k(0)$ for $k\le3$ the midpoint coefficients and $c_4(\xi)$ a fourth
  coefficient at some interior point. Integrating $\sigma^k$ over $[-1,1]$ kills
  the odd powers and gives $\int_{-1}^1\sigma^2=\tfrac23$,
  $\int_{-1}^1\sigma^4=\tfrac25$; multiplying by $h=\tfrac{r-l}2$ yields the
  stated combination $(r-l)(c_0+\tfrac{c_2}3+\tfrac{c_4}5)$, and the whole-panel
  jet supplies the uniform upper bound for $c_4$.
\end{proof}
\formal{panelUpper, integral\_le\_panelUpper, affineJet, centerJet}{Upper0423/Native/Majorant.lean, Upper0423/Native/PanelSound.lean}

The rule is a genuine upper bound for the integral, not an approximation. The
fourth-order Taylor form is what makes the certificates small: on a smooth panel
it converges like $(r-l)^5$ while the range rule converges only like $(r-l)$.

\subsection{Panel majorants}
\label{sec:modulus}

A certificate leaf fixes a \emph{box}
\[
  b=\bigl(n,\ \beta_{\mathrm{lo}},\beta_{\mathrm{hi}},\
  \eta_{\mathrm{lo}},\eta_{\mathrm{hi}},\ c,\ \tau,\ \varepsilon\bigr)
\]
of rational data, subject to the validity conditions
\[
  n\ge2,\quad
  1\le\beta_{\mathrm{lo}}\le\beta_{\mathrm{hi}},\quad
  0\le\eta_{\mathrm{lo}}\le\eta_{\mathrm{hi}}\le1,\quad
  c>0,\quad 0<\varepsilon<\tau\le\tfrac12,\quad
  c\tau<\tfrac32,\quad c<\pi .
\]
The derived quantities used by the majorants are
\begin{gather*}
  \delta=\beta_{\mathrm{hi}}-1,\qquad
  d=\min\Bigl(1-\tfrac1{\beta_{\mathrm{hi}}},\,1-\eta_{\mathrm{lo}}\Bigr),\qquad
  a=\beta_{\mathrm{hi}}+\eta_{\mathrm{hi}},\\
  \Sigma_1=\delta\bigl(\beta_{\mathrm{hi}}+\tfrac53\bigr),\qquad
  \Sigma_2=\beta_{\mathrm{hi}}^2-\eta_{\mathrm{lo}}^2,\qquad
  J=\tfrac{1}{36}\min(\Sigma_1,\Sigma_2) .
\end{gather*}
The reason $d$ is a legitimate upper bound for $1-s$ is the pair of inequalities
$s\ge1/\beta$ (from $s\beta\ge1$, Lemma~\ref{lem:eta}) and $s\ge\eta$; the
quantities $\Sigma_1,\Sigma_2$ dominate $(\beta-1)(\beta+5/3)$ and
$\beta^2-\eta^2$ respectively.
\formal{Box, Box.valid, Box.dcap, Box.dsum, Box.cubic1, Box.cubic2, Box.imagSq, dcap\_ge, dsum\_ge, cubic1\_ge, cubic2\_ge}{Upper0423/Native/Majorant.lean, Upper0423/Analytic/ModulusBridge.lean}

With $u=ct$, the four \emph{branch choices} of a panel are as follows.

\paragraph{Imaginary majorant $j(u)$.} One of
\[
  \tfrac16\sqrt{\Sigma_1}\,u^3,\quad
  \tfrac12\sqrt{\tfrac23\delta}\,u^2,\quad
  \tfrac16\sqrt{\Sigma_2}\,u^3,\quad
  \sqrt{2d}\,(\sin u-u\cos u)+d\,u^2,
\]
dominating $|\operatorname{Im}f(u)|$ by Propositions~\ref{prop:imcubic},
\ref{prop:imquad}, \ref{prop:circle}, \ref{prop:radial} respectively.

\paragraph{Real majorant $\mathrm{re}(u)$.} One of
\[
  \cos u+\tfrac16\delta u^3,\quad
  \cos u+d\,(u\sin u+u^2),\quad
  1-\tfrac{u^2}2+\kap\beta_{\mathrm{hi}}u^3,
\]
dominating $\operatorname{Re}f(u)$ by Proposition~\ref{prop:reupper}.

\paragraph{Modulus majorant $M(u)$.} One of
\begin{align*}
  \text{\emph{one}:}\quad& 1;\\
  \text{\emph{line}:}\quad& \sqrt{1-u^2+2\kap a\,u^3};\\
  \text{\emph{cosine}:}\quad& \sqrt{1-\dfrac{2A\bigl(1-\cos(a u)\bigr)}{a^2}};\\
  \text{\emph{coordinate}:}\quad& \sqrt{\mathrm{far}(\cos u,\mathrm{re}(u);G)+j(u)^2} .
\end{align*}
The first three are $|f|\le1$ and Propositions~\ref{prop:rekappa},
\ref{prop:profilebound}; the \emph{cosine} branch is guarded by a whole-panel
check that $4\le a\,c\,l$ and $a\,c\,r\le2\pi$, so that the curved branch of the
profile applies on the entire panel. The \emph{coordinate} branch is the
disk/strip intersection: $|f|^2=(\operatorname{Re}f)^2+(\operatorname{Im}f)^2$,
the imaginary square is bounded by $j(u)^2$, and the real square is bounded by
\[
  \mathrm{far}(\alpha,\gamma;G)=
  \begin{cases}
    \gamma^2, & G=+1,\\
    \alpha^2, & G=-1,\\
    \alpha^2+\gamma^2, & G=0,
  \end{cases}
\]
where $\alpha\le\operatorname{Re}f(u)\le\gamma$ are the lower bound $\cos u$
(Proposition~\ref{prop:relower}) and the chosen upper bound, and $G$ is a
\emph{sign guard} computed once on the whole panel: $G=+1$ if $\alpha+\gamma\ge0$
throughout the panel (so $|\operatorname{Re}f|\le\gamma$), $G=-1$ if
$\alpha+\gamma\le0$ throughout (so $|\operatorname{Re}f|\le-\alpha$), and $G=0$
otherwise. Since $\max(|\alpha|,|\gamma|)^2\le\alpha^2+\gamma^2$ always, the
guard only improves the bound, and its correctness on the whole panel is what
allows the sharper midpoint jet to reuse it. The idea of intersecting the Taylor
disk with the real and imaginary strips, which is what this branch implements, is
due to [XL] (item (vi) of Section~\ref{sec:attribution}); the sign guard and the
use of the two-sided real bracket of Proposition~\ref{prop:relower} are
particular to the present engine.
\formal{imagJet, realJet, modulusJet, cosineGuard, signGuard, farthestSquare, abs\_im\_le\_imagR, re\_le\_realR, norm\_le\_modulusR, cosineGuard\_sound}{Upper0423/Native/Majorant.lean, Upper0423/Analytic/ModulusBridge.lean, Upper0423/Analytic/IntegrandBridge.lean}

\subsection{The low-frequency error factor}
\label{sec:lowfactor}

The low-frequency integrand \eqref{eq:integrands} needs a bound for
$|f(u)^n-g(u)^n|$, and this is where the geometric-sum comparison enters. Writing
$G_n(y,z)=\sum_{j<n}y^jz^{n-1-j}$, the telescoping identity gives
\begin{equation}
  \label{eq:telescope}
  \bigl|f(u)^n-g(u)^n\bigr|\le\bigl|f(u)-g(u)\bigr|\cdot G_n\bigl(M(u),g(u)\bigr)
\end{equation}
whenever $|f|\le M$ and $g\ge0$. The engine computes $G_n$ exactly by the
halving recursion
\[
  G_{2k}(y,z)=(y^k+z^k)\,G_k(y,z),\qquad
  G_{2k+1}(y,z)=y\,G_{2k}(y,z)+z^{2k},
\]
returning $(y^n,z^n,G_n)$ from $(y^{\lfloor n/2\rfloor},z^{\lfloor n/2\rfloor},
G_{\lfloor n/2\rfloor})$; this costs $O(\log n)$ jet multiplications.
\formal{geometric, geometricSum, geomSum, geomSum\_double, geomSum\_succ\_double, geometricR\_spec, norm\_pow\_sub\_pow\_le\_geomSum}{Upper0423/Native/Majorant.lean, Upper0423/Analytic/Geometric.lean}

Five branches for the factor $\Lambda(u)$ with
$|f(u)^n-g(u)^n|\le u^3\Lambda(u)$ (or the trivial split) are provided:

\begin{itemize}[leftmargin=1.6em]
\item \emph{classical}: $\Lambda=\bigl(\tfrac16\beta_{\mathrm{hi}}
  + \mathrm{GU}(u)\bigr)\,G_n(M,g)$, where
  $g(u)-1+\tfrac{u^2}2\le u^3\,\mathrm{GU}(u)$ and
  $\|f(u)-1+\tfrac{u^2}2\|\le\tfrac16\beta|u|^3$ by \eqref{eq:quadratic}: the two
  Taylor disks around the common center $1-\tfrac{u^2}2$.
\item \emph{coordinate}: the real and imaginary parts of $f(u)-g(u)$ are
  bracketed separately. Using the two-sided Taylor brackets
  $u^3\,\mathrm{GC}_L(u)\le g(u)-\cos u\le u^3\,\mathrm{GC}_U(u)$ one gets
  $\alpha'\le\operatorname{Re}(f(u)-g(u))\le\gamma'$ with
  $\alpha'=-\mathrm{GC}_U(u)u^3$ and
  $\gamma'=\bigl(\tfrac{\delta}{6}-\mathrm{GC}_L(u)\bigr)u^3$, so that
  $\Lambda=\sqrt{\mathrm{far}(\alpha',\gamma';G')+J}\cdot G_n(M,g)$ with a second
  sign guard $G'$ and the imaginary square cap $J$.
\item \emph{anchor}: split around $\cos u$:
  $\Lambda=\sqrt{\tfrac{\delta^2}{36}+J}\;G_n(M,\cos u)
  + \mathrm{GC}_U(u)\,G_n(g,\cos u)$.
\item \emph{ode}: the same split with the ODE bound of
  Proposition~\ref{prop:ode} in place of the anchor disk:
  $\Lambda=\sqrt{\Sigma_1}\,\mathrm{SU}(u)\,G_n(M,\cos u)
  + \mathrm{GC}_U(u)\,G_n(g,\cos u)$, where $u-\sin u\le u^3\mathrm{SU}(u)$.
\item \emph{trivial}: $|f^n-g^n|\le M^n+g^n$, used only where the difference
  bound is worthless.
\end{itemize}
The polynomials $\mathrm{GU},\mathrm{GC}_L,\mathrm{GC}_U,\mathrm{SU}$ are the
explicit alternating Taylor truncations evaluated by the engine, and the
corresponding real inequalities are proved on $0\le u\le\pi/2$ in
\file{Analytic/TaylorPolys.lean} from Mathlib's exponential and trigonometric
Taylor bounds.
\formal{lowFactor, lowerErrorJet, upperErrorJet, gaussianCosineCubic, sineCubicUpper, gaussianCubicUpper}{Upper0423/Native/Majorant.lean}
\formal{gaussianCubicUpper\_bound, gaussianCosine\_upper\_bound, gaussianCosine\_lower\_bound, sineCubicUpper\_bound, norm\_pow\_sub\_gauss\_le\_lowFactor, low\_integrand\_le, high\_integrand\_le, normal\_integrand\_le}{Upper0423/Analytic/TaylorPolys.lean, Upper0423/Analytic/IntegrandBridge.lean}

\subsection{The panel integrand and its certified bound}

Putting the pieces together, the engine's real majorant for a panel $[l,r]$ of
part $p\in\{\text{low},\text{high},\text{normal}\}$ with branch choices and
guards is
\[
  \mathrm{build}_b(t)=
  \begin{cases}
    |K(t)|_{\text{normal}}\cdot e^{-n c^2t^2/2}, & p=\text{normal},\\[2pt]
    |K(t)|_{\text{left/right}}\cdot M(ct)^n, & p=\text{high},\\[2pt]
    c^3t^2\cdot\bigl(s|K(s)|\bigr)\big|_{s=t}\cdot\Lambda(ct), & p=\text{low},\\[2pt]
    |K(t)|_{\text{left}}\cdot\bigl(M(ct)^n+e^{-c^2t^2/2\cdot n}\bigr),
      & p=\text{low, trivial}.
  \end{cases}
\]
(The low case uses form \eqref{eq:K-left}: $c^3t^2\cdot s|K(s)|=u^3|K(t)|$ with
$u=ct$.) Comparing with \eqref{eq:integrands}, each case dominates the
corresponding Prawitz integrand pointwise on the panel, by the results of
Sections~\ref{sec:chf}, \ref{sec:modulus} and \ref{sec:lowfactor}. The function
\Lean{framePanel} then

\begin{enumerate}[label=(\arabic*),leftmargin=2.2em]
\item checks the box validity, the panel validity ($\varepsilon\le l<r\le\tau$ for
  low, $\tau\le l<r\le1$ and no straddling of $\tfrac12$ for high,
  $0\le l<r\le\tau$ for normal) and the branch validity;
\item evaluates the jet of $\mathrm{build}_b$ on the whole panel, \emph{computing}
  the sign guards from it;
\item evaluates the jet at the midpoint, \emph{reusing} the stored guards;
\item returns the panel rule value of Lemma~\ref{lem:panel}, or fails.
\end{enumerate}

\begin{theorem}[Panel soundness]
  \label{thm:panelsound}
  If $\Lean{framePanel}\,b\,l\,r\,p\,\mathrm{ch}=\mathrm{some}\,U$, then $b$ and
  the panel are valid, $U\ge0$, and there are guards $g$ valid on all of
  $[l,r]$ with
  \[
    \int_l^r\mathrm{build}_b(t)\dd t\le U .
  \]
  If moreover $\mu$ is admissible with $(\beta,\eta)$ in the box, then
  \[
    \int_l^r(\text{the corresponding Prawitz integrand})\dd t\le U .
  \]
\end{theorem}
\formal{framePanel\_sound, framePanel\_sound', lowPanel\_integral\_le, highPanel\_integral\_le, normalPanel\_integral\_le}{Upper0423/Native/PanelSound.lean, Upper0423/Small/Chains.lean, Upper0423/Small/Leaf.lean}

\subsection{Chains, omission and tail}

A \emph{chain} is a list of adjacent panels covering an interval; the chain total
bounds the integral over the whole interval, provided the endpoints match, which
the checker verifies explicitly (the low chain must end exactly at $\tau$, the
high chain at $1$, the normal chain at $\tau$).
\formal{chainSum, chainSum\_integral\_le}{Upper0423/Native/Cover.lean, Upper0423/Small/Chains.lean}

The interval $[0,\varepsilon]$ of the low part is not covered by panels: the
kernel is singular at $0$ and the Taylor jets degenerate. Instead it is bounded
in closed form.

\begin{lemma}[Low-frequency omission]
  \label{lem:omission}
  Let $0<\varepsilon\le\tfrac12$, $c>0$, $n\ge1$ and $\beta\le\beta_{\mathrm{hi}}$.
  Then
  \[
    \int_0^\varepsilon|K(t)|\,\bigl|f(ct)^n-g(ct)^n\bigr|\dd t
    \le n\Bigl[\beta_{\mathrm{hi}}c^3\Bigl(\frac{\varepsilon^4}{24}
      +\frac{\varepsilon^3}{18\pi}\Bigr)
      + c^4\Bigl(\frac{\varepsilon^5}{40}+\frac{\varepsilon^4}{32\pi}\Bigr)\Bigr] .
  \]
\end{lemma}
\begin{proof}[Proof sketch]
  For $0<t\le\tfrac12$ one has $\cot(\pi t)\le\tfrac1{\pi t}$, whence
  $0\le k(t)\le\tfrac1{\pi t}$ and $|K(t)|\le1+\tfrac1{\pi t}$. Also
  $|f^n-g^n|\le n|f-g|\le n\bigl(\tfrac{\beta u^3}{6}+\tfrac{u^4}{8}\bigr)$ with
  $u=ct$, by the two Taylor disks around $1-\tfrac{u^2}2$. Multiplying and
  integrating $t^2,t^3,t^4$ over $[0,\varepsilon]$ gives the stated expression.
\end{proof}
\formal{norm\_prawitzKernel\_le\_one\_add, omission\_bound, omissionUpper, omissionUpper\_sound}{Upper0423/Small/Omission.lean, Upper0423/Native/Cover.lean}

The Gaussian tail \eqref{eq:tail} is handled by panelling
$v\mapsto e^{-v}/v$ on $[x,x+24]$ with deterministic cut points (doubling from
$x$ until $\tfrac14$, then steps of $\tfrac14$), plus the closed tail bound
$\int_R^\infty e^{-v}/v\dd v\le e^{-R}/R$.
\formal{tailCuts, tailChain, tailUpper, tailPanel\_sound, tailChain\_sound, integral\_exp\_neg\_div\_Ioi\_le, tailUpper\_sound}{Upper0423/Native/Cover.lean, Upper0423/Small/Tail.lean, Upper0423/Native/PanelSound.lean}

\subsection{A leaf and its budget}

\begin{definition}
  \label{def:leaf}
  Given a box $b$ and leaf data $d$ (the cutoff $c$, the split $\tau$ and the
  three chains), the \emph{leaf total} is
  \[
    U(b,d)=\underbrace{\Sigma_{\text{low}}}_{[\varepsilon,\tau]}
    +\underbrace{\Sigma_{\text{high}}}_{[\tau,1]}
    +\underbrace{\Sigma_{\text{normal}}}_{[0,\tau]}
    +\underbrace{\Omega}_{[0,\varepsilon]}
    +\underbrace{\mathcal T}_{\text{tail}} ,
  \]
  and the leaf is \emph{accepted} when $b$ is valid, $U\ge0$ and
  \begin{equation}
    \label{eq:budget}
    n\,U^2 \;<\; \Bigl(\frac{423}{1000}\Bigr)^2\beta_{\mathrm{lo}}^2 .
  \end{equation}
\end{definition}
\formal{leafTotal, leafAccepted}{Upper0423/Native/Cover.lean}

\begin{theorem}[Leaf soundness]
  \label{thm:leaf}
  Let $\mu$ be admissible with $\beta_{\mathrm{lo}}\le\beta\le\beta_{\mathrm{hi}}$
  and $\eta_{\mathrm{lo}}\le\eta\le\eta_{\mathrm{hi}}$, and suppose the leaf
  $(b,d)$ is accepted. Then for every $x$,
  \[
    \mathcal E(\mu,n,x)\le\frac{423}{1000} .
  \]
\end{theorem}
\begin{proof}
  Theorem~\ref{thm:prawitz} with $T=c\sqrt n$ and split $\tau$ bounds
  $|F_n(x)-\Phi(x)|$ by the smoothing integral; splitting the low integral at
  $\varepsilon$ and applying the chain, omission and tail bounds gives
  $|F_n(x)-\Phi(x)|\le U$. The budget \eqref{eq:budget} is equivalent to
  $\sqrt n\,U<\tfrac{423}{1000}\beta_{\mathrm{lo}}$ (both sides nonnegative),
  hence
  \[
    \mathcal E(\mu,n,x)=\frac{\sqrt n}{\beta}|F_n-\Phi|
    \le\frac{\sqrt n\,U}{\beta}
    \le\frac{423}{1000}\cdot\frac{\beta_{\mathrm{lo}}}{\beta}
    \le\frac{423}{1000},
  \]
  using $\beta_{\mathrm{lo}}\le\beta$.
\end{proof}
\formal{leafTotal\_cdf\_bound, leafAccepted\_normalizedError}{Upper0423/Small/Leaf.lean}

Note the shape of \eqref{eq:budget}: it is a purely rational inequality in the
leaf data, which is what makes it decidable, and the square is taken to avoid
$\sqrt n$ inside the checker.

\section{The per-sample-size covers, $2\le n\le255$}
\label{sec:covers}

\subsection{Cover trees}

\begin{definition}
  \label{def:cover}
  A \emph{cover} is a binary tree whose leaves carry leaf data and whose internal
  nodes carry an axis and a rational split point. Acceptance over a root
  rectangle is defined recursively: a leaf is accepted if the leaf box obtained
  by outward rounding of the rectangle to the $2^{-40}$ grid is accepted in the
  sense of Definition~\ref{def:leaf}; a split node is accepted if the split point
  lies strictly inside the rectangle along the chosen axis and both children are
  accepted over the two halves.
\end{definition}
\formal{Cover, coverAccepted, coverAcceptedBE, roundDown, roundUp, imageBox, imageBoxBE}{Upper0423/Native/Cover.lean, Upper0423/Native/CoverGen.lean}

\begin{lemma}[Locating a law in a cover]
  \label{lem:locate}
  If a cover is accepted over a rectangle $R$ and $(\beta,\eta)$ (respectively
  $(\delta,w)$) lies in $R$, then some accepted leaf box contains
  $(\beta,\eta)$.
\end{lemma}
\begin{proof}
  Induction on the tree: at a split node, one of the two halves contains the
  point because the split point lies inside the rectangle; at a leaf, outward
  rounding only enlarges the box.
\end{proof}
\formal{coverAccepted\_leaf, cov\_coverAcceptedBE\_leaf, inBox\_imageBox, cov\_inBox\_imageBoxBE}{Upper0423/Small/CoverTree.lean, Upper0423/Covers/Sound.lean}

Two coordinate systems are used, as anticipated in Section~\ref{sec:moments}.
For $\beta\le2$, the cover is a tree in $(\delta,w)\in[\delta_{\mathrm{lo}},
\delta_{\mathrm{hi}}]\times[0,1]$ with $\beta=1+\delta$ and $\eta=1-\delta w$;
the map $(\delta,w)\mapsto(\beta,\eta)$ is well defined because, given
$\delta\ge0$ and $\eta\in[\max(0,1-\delta),1]$, the value
$w=(1-\eta)/\delta$ (or $w=0$ if $\delta=0$) lies in $[0,1]$ and satisfies
$1-\delta w=\eta$; here $\eta\ge1-\delta$ is exactly
Proposition~\ref{prop:feasible}. For $\beta\ge2$ the cover is a tree directly in
$(\beta,\eta)\in[\beta_{\mathrm{lo}},\beta_{\mathrm{hi}}]\times[0,1]$.

\subsection{The certificate format}

Each cover is serialized as a whitespace-separated token string and parsed by an
untrusted recursive-descent parser: \verb|S axis m <left> <right>| for a split
and
\[
  \verb|F c tau nL (r code)* nH (r code)* nN r*|
\]
for a leaf, where \verb|code| packs the four branch choices into a single natural
number ($\mathrm{imag}+4\,\mathrm{real}+12\,\mathrm{modulus}
+48\,\mathrm{error}$). Parsing is untrusted precisely because the soundness
theorem is universally quantified over \emph{every} tree the checker accepts: if
the parser produced garbage, the checker would simply reject it. A representative
leaf, from \file{Covers/Data/BE3.lean}, reads
\begin{leancode}
F 477/200 2/5 17 1/500 216 1/250 216 ... 20 17/40 20 ... 13 2/65 4/65 ... 2/5
\end{leancode}
declaring $c=477/200$, $\tau=2/5$, $17$ low panels, then the high panels and the
$13$ normal panel endpoints.
\formal{parseCover, tokens, coverStringAccepted, coverStringAcceptedDW, coverStringAcceptedBE}{Upper0423/Native/Cover.lean, Upper0423/Native/CoverGen.lean}

\subsection{The certificates}

Table~\ref{tab:covers} summarizes the finite data. The counts were obtained by
direct inspection of the repository.

\begin{table}[htbp]
  \centering
  \begin{tabular}{lrrl}
    \toprule
    Family & Files & Leaves & Root domain\\
    \midrule
    \file{Small/Data/Cover\{2..63\}}  & $62$  & $5\,053$  & $\delta\in[0,\tfrac{258569}{454000}]$, $w\in[0,1]$\\
    \file{Covers/Data/DW\{2..255\}}   & $254$ & $8\,953$  & $\delta\in[\delta_n,d_n]$, $w\in[0,1]$\\
    \file{Covers/Data/BE\{3..255\}}   & $253$ & $4\,839$  & $\beta\in[2,\beta_n]$, $\eta\in[0,1]$\\
    \midrule
    total                             & $569$ & $18\,845$ & \\
    \bottomrule
  \end{tabular}
  \caption{The per-sample-size certificates. Here $\beta_n$ is a rational
    number with $\beta_n\ge\tfrac{13003}{10000}\sqrt n$;
    $\delta_n=\tfrac{258569}{454000}$ for $n\le63$ and $\delta_n=0$ for
    $n\ge64$; and $d_n=\min(1,\beta_n-1)$ (so $d_n=1$ except
    $d_2=\tfrac{17}{20}$). Each file carries exactly one
    \Lean{native\_decide}.}
  \label{tab:covers}
\end{table}

For each $n$ with $2\le n\le255$ the file \file{Covers/All.lean} contains a
theorem \Lean{cover\_}$n$ of the form
\[
  \beta\le\beta_n\ \Longrightarrow\
  \mathcal E(\mu,n,x)\le\tfrac{423}{1000},
\]
proved by the case distinction
\begin{itemize}[leftmargin=1.6em]
\item $\beta\le\tfrac{712569}{454000}$ and $n\le63$: the original research cover
  \file{Small/Data/Cover}$n$;
\item $\tfrac{712569}{454000}<\beta\le2$ (or $\beta\le2$ when $n\ge64$): the
  cover \file{Covers/Data/DW}$n$;
\item $\beta>2$: the cover \file{Covers/Data/BE}$n$ (absent for $n=2$, where
  $\beta\le\beta_2=37/20<2$ makes the branch vacuous).
\end{itemize}
The hypothesis $\beta\le\beta_n$ is supplied from
$\ell=\beta/\sqrt n<\tfrac{13003}{10000}$ by the elementary
\Lean{beta\_le\_of\_ell}: if $b\ge0$ and $(\tfrac{13003}{10000})^2n\le b^2$ then
$\sqrt n\cdot\tfrac{13003}{10000}\le b$, so $\beta<b$. The rational $\beta_n$ is
the smallest convenient rational above $\tfrac{13003}{10000}\sqrt n$; e.g.
$\beta_2=37/20$, $\beta_3=227/100$, $\beta_{64}=521/50$,
$\beta_{255}=1039/50$. Finally \Lean{normalizedError\_moderate\_le} does
\Lean{interval\_cases n} and applies \Lean{cover\_}$n$.
\formal{beta\_le\_of\_ell, cover\_2 \dots cover\_255, normalizedError\_moderate\_le}{Upper0423/Covers/All.lean}
\formal{normalizedError\_le\_of\_coverAccepted, normalizedError\_small\_le}{Upper0423/Small/Leaf.lean, Upper0423/Small/All.lean}
\formal{normalizedError\_le\_of\_dwCover, normalizedError\_le\_of\_beCover}{Upper0423/Covers/Sound.lean}

\begin{remark}
  The three families overlap deliberately: the \file{Small} covers were produced
  first, for the narrow range $\beta\le B=712569/454000$, and were retained
  rather than regenerated. The reason for the separate treatment of $\beta\le2$
  and $\beta\ge2$ is that in the first regime $\eta$ is pinned near $1$ by
  $\beta+\eta\ge2$, so the $(\delta,w)$ coordinates are much better adapted; for
  $\beta\ge2$ the constraint is inactive and $(\beta,\eta)$ is the natural chart.
\end{remark}

\section{The uniform certificate, $n\ge256$}
\label{sec:uniform}

For large $n$, a per-$n$ cover is impossible. Instead the whole smoothing
functional is reduced to five one-dimensional integrals in two scalar parameters,
and a single certificate covers all $n\ge256$ simultaneously.

\subsection{The two scalars}

Fix the smoothing parameters
\begin{equation}
  \label{eq:uniform-params}
  a=\beta+\eta,\qquad
  L=\frac{a}{\sqrt n},\qquad
  h=\frac1{\sqrt n},\qquad
  T=\frac{2\pi}{L}=\frac{2\pi\sqrt n}{a},\qquad
  \tau=\frac14 .
\end{equation}
With this choice, \eqref{eq:Tscale} holds with $c=2\pi/a$, so the per-summand
frequency \eqref{eq:u} is $u=2\pi t/a$, and
\[
  a\,u=2\pi t,\qquad n\,u^2=\frac{(2\pi t)^2}{L^2},\qquad
  (n-1)\,u^2=\frac{(1-h^2)(2\pi t)^2}{L^2} .
\]
Thus every occurrence of $n$, $\beta$ and $\eta$ in the integrands collapses into
$L$ and $h$.

\begin{definition}
  \label{def:uniform}
  With $q$ the convex profile of Definition~\ref{def:profile} and
  $\psi(t)=(2\pi t)^2q(2\pi t)$, set
  \begin{align*}
    \mathrm{high}(L)&=\int_{1/4}^{1}|K(t)|\,e^{-\psi(t)/L^2}\dd t,\\
    \mathrm{norm}(L)&=\int_0^{1/4}\Bigl|K(t)-\frac{\ii}{\pi t}\Bigr|
      e^{-2\pi^2t^2/L^2}\dd t,\qquad
    \mathrm{tail}(L)=\int_{1/4}^\infty\frac{e^{-2\pi^2t^2/L^2}}{\pi t}\dd t,\\
    E_1(L,h,t)&=\exp\Bigl(-\frac{(1-h^2)(2\pi t)^2}{2L^2}\Bigr),\qquad
    E_2(L,h,t)=\exp\Bigl(-\frac{(1-h^2)(2\pi t)^2
      \bigl(\tfrac12-2\pi\kap t\bigr)}{L^2}\Bigr),\\
    \mathrm{lowG}(L,h)&=\int_0^{1/4}|K(t)|\,\frac{(2\pi t)^3}{12L^3}
      \bigl(E_1+E_2\bigr)\dd t,\qquad
    \mathrm{lowK}(L,h)=\int_0^{1/4}|K(t)|\,\frac{(2\pi t)^4}{12L^4}E_1\dd t .
  \end{align*}
\end{definition}
\formal{psiR, highF, normalF, tailF, lowE1, lowE2, lowG, lowK}{Upper0423/Uniform/Defs.lean}

Note $E_1=g(u)^{n-1}$ and $E_2=m(u)^{n-1}$ with $m(u)=e^{-u^2/2+\kap a u^3}$, so
these are exactly the two exponential envelopes the geometric sum produces.

\begin{proposition}[Uniform reduction]
  \label{prop:reduction}
  For every admissible $\mu$, every $n\ge1$ and every $x$,
  \begin{equation}
    \label{eq:reduction}
    \mathcal E(\mu,n,x)\;\le\;
    \frac{a}{\beta}\cdot
      \frac{\mathrm{high}(L)+\mathrm{norm}(L)+\mathrm{tail}(L)}{L}
    +\frac{\Gamma}{\beta}\,\mathrm{lowG}(L,h)
    +\frac{h}{\beta}\,\mathrm{lowK}(L,h) ,
  \end{equation}
  with $a,L,h$ as in \eqref{eq:uniform-params} and $\Gamma$ as in
  \eqref{eq:gamma}.
\end{proposition}
\begin{proof}[Proof sketch]
  Apply Theorem~\ref{thm:prawitz} with $T=2\pi/L$ and $\tau=\tfrac14$ and bound
  the four terms.

  \emph{High.} By Proposition~\ref{prop:profilebound},
  $|f(u)|\le e^{-u^2q(a|u|)}$, hence
  $|f(u)|^n\le e^{-nu^2q(2\pi t)}=e^{-\psi(t)/L^2}$; the deduction
  $|f|\le e^{-u^2q}$ from $|f|^2\le1-2u^2q$ uses $1-2y\le e^{-2y}$.

  \emph{Low.} By the geometric-sum comparison and the pointwise estimates of
  Propositions~\ref{prop:ode}, \ref{prop:rekappa} one has, for
  $0\le u\le\pi/2$,
  \[
    \bigl|f(u)^n-g(u)^n\bigr|
    \le \frac{n\Gamma u^3}{12}\bigl(g(u)^{n-1}+m(u)^{n-1}\bigr)
      +\frac{n u^4}{12}\,g(u)^{n-1} ,
  \]
  where $\Gamma=\sqrt{(\beta-1)^2+\beta^2-\eta^2}$ arises as the radius of the
  disk containing $f(u)-\cos u$ jointly with the imaginary strip. The constraint
  $u\le\pi/2$ holds on $t\le\tfrac14$ because $u=2\pi t/a\le\pi/(2a)\le\pi/4$ by
  Proposition~\ref{prop:feasible}. Multiplying by $|K(t)|$, integrating, and
  substituting $u=2\pi t/a$ gives
  $\tfrac{\Gamma}{\sqrt n}\mathrm{lowG}+\tfrac1n\mathrm{lowK}$.

  \emph{Normal and tail.} These do not involve $\mu$ at all and are literally
  $\mathrm{norm}(L)$ and $\mathrm{tail}(L)$.

  Finally $\tfrac{\sqrt n}{\beta}=\tfrac{a}{\beta}\cdot\tfrac1L$ and
  $\tfrac{\sqrt n}{\beta}\cdot\tfrac{1}{n}=\tfrac{h}{\beta}$, which produces
  \eqref{eq:reduction}.
\end{proof}
\formal{red\_high\_pointwise, red\_low\_pointwise, red\_low\_integral, red\_high\_integral, normalized\_cdf\_le\_uniform}{Upper0423/Uniform/Reduction.lean}
\formal{norm\_charFun\_sub\_cos\_le\_gamma, norm\_charFun\_le\_cubicExp, low\_geomSum\_le, norm\_charFun\_pow\_sub\_gauss\_le}{Upper0423/LargeN/LowPointwise.lean}

\subsection{The budget}

Let $z=\eta/\beta\in[0,1]$, so $a/\beta=1+z$. Suppose $F$, $G$, $K$ are rational
upper bounds for the three bracketed quantities in \eqref{eq:reduction} on a cell
$L\in[L_1,L_2]$. Two constraints convert \eqref{eq:reduction} into a rational
inequality.

\begin{lemma}[$\Gamma$ constraint]
  \label{lem:gammaconstraint}
  If $\beta\ge1$ and $0\le\eta\le1$ then
  $\bigl(\Gamma/\beta\bigr)^2\le2(1-\eta/\beta)$.
\end{lemma}
\begin{proof}
  Equivalently $(\beta-1)^2+\beta^2-\eta^2\le2\beta^2-2\beta\eta$, i.e.
  $(1-\eta)(2\beta-1-\eta)\ge0$, which holds since $\eta\le1$ and
  $2\beta-1-\eta\ge2-1-1=0$.
\end{proof}

\begin{lemma}[$z$ and $h$ caps]
  \label{lem:caps}
  Let $r\ge2$ with $r^2\le n$, so $hr\le1$. On a cell with
  $L_1\le(\beta+\eta)h\le L_2$:
  \[
    z\le z_{\mathrm{hi}}:=
    \begin{cases}
      \min\bigl(1,\tfrac{1}{L_1r-1}\bigr), & L_1r>1,\\
      1, & \text{otherwise},
    \end{cases}
    \qquad
    \frac h\beta\le \min\Bigl(\frac1r,\frac{L_2}{2},
      \frac{1+z_{\mathrm{hi}}}{r^2L_1}\Bigr) .
  \]
\end{lemma}
\begin{proof}
  From $L_1\le(\beta+\eta)h$ and $hr\le1$ we get
  $L_1r\le(\beta+\eta)hr\le\beta+\eta$, hence $\beta\ge L_1r-1$ (as $\eta\le1$)
  and $z=\eta/\beta\le1/(L_1r-1)$. For the $h$ caps: $h\le1/r$; and
  $2h\le(\beta+\eta)h\le L_2$ by Proposition~\ref{prop:feasible};
  and $h\,r^2L_1\le(hr)^2(\beta+\eta)\le\beta+\eta\le(1+z_{\mathrm{hi}})\beta$.
\end{proof}

\begin{lemma}[Budget algebra]
  \label{lem:budget}
  Let $c>0$, $0\le z\le z_{\mathrm{hi}}$, $0\le v$ with $v^2\le2(1-z)$,
  $0\le k\le k_c$, $K\ge0$, and $F'\le F$, $G'\le G$, $K'\le K$. Then
  \[
    (1+z)F'+vG'+kK'
    \;\le\; F+c+\frac{G^2}{2c}+\max\bigl(0,z_{\mathrm{hi}}(F-c)\bigr)+k_cK .
  \]
\end{lemma}
\begin{proof}
  $vG\le \tfrac{c v^2}{2}+\tfrac{G^2}{2c}$ (AM--GM, i.e.
  $(cv-G)^2\ge0$) and $\tfrac{cv^2}{2}\le c(1-z)$, so
  \[
    (1+z)F+vG\le F+zF+c-cz+\frac{G^2}{2c}
    = F+c+z(F-c)+\frac{G^2}{2c},
  \]
  and $z(F-c)\le\max(0,z_{\mathrm{hi}}(F-c))$ in both signs of $F-c$.
\end{proof}
\formal{bud\_gamma\_sq, bud\_z\_le, bud\_hcap, bud\_kc, bud\_k\_le, bud\_core, bud\_val, bud\_accept, bud\_accept\_sound}{Upper0423/Uniform/Budget.lean}

A cell $[L_1,L_2]$ is therefore \emph{accepted} when $2\le r$,
$0<L_1\le L_2\le3$, the three certified bounds $F=U_1+U_2$, $G$, $K$ exist, and
\[
  F+c+\frac{G^2}{2c}+\max\bigl(0,z_{\mathrm{hi}}(F-c)\bigr)+k_cK
  \;<\;\frac{423}{1000}
\]
for at least one of the three candidate parameters
$c\in\bigl\{F,\ G/\sqrt2,\ G/\sqrt{2(1-z_{\mathrm{hi}})}\bigr\}$
(the square roots replaced by rational under-approximations; the choice of
candidate is immaterial to soundness, since \emph{any} $c>0$ works in
Lemma~\ref{lem:budget}).

\subsection{Certified bounds for the five integrals}

It remains to produce $U_1\ge(\mathrm{high}+\mathrm{tail})/L$,
$U_2\ge\mathrm{norm}/L$, $G\ge\mathrm{lowG}$ and $K\ge\mathrm{lowK}$, uniformly
for $L\in[L_1,L_2]$ and $h\le h_{\mathrm{cap}}$. All four are monotone in the
right places, which is what makes a cell-uniform bound possible.

\paragraph{High and tail.} Both integrands increase with $L$, so they are
evaluated at $L_2$. For the high part, $\psi$ is bounded below by one of two
explicit branches: the supporting line
$\psi(t)\ge(2\pi t)^2(\tfrac12-\kap\,2\pi t)$ (valid for $2\pi t\le2\pi$), and
the exact curved branch $\psi(t)=A(1-\cos2\pi t)$ once $2\pi t\ge4$. The kernel is
bounded by \eqref{eq:K-left} or \eqref{eq:K-right} according to the panel. The
resulting smooth majorant is integrated by the panel rule of
Lemma~\ref{lem:panel} over $\lfloor m/3\rfloor+1$ panels of $[\tfrac14,1]$ and
the tail over the remaining panels of $[\tfrac14,S]$, plus a closed Gaussian tail
beyond $S$.
\formal{ht\_psiJet, ht\_highPanel, ht\_tailPanel, ht\_tailRest, highTailUpper, highTailUpper\_sound, highTail\_small}{Upper0423/Uniform/HighTail.lean}

\paragraph{Normal correction.} Two bounds are computed and the smaller kept. The
first is closed-form, from the quadratic kernel majorant
$|K(t)-\ii/(\pi t)|\le p(t)=1-t+\nu t^2$ on $(0,\tfrac14]$ (with $\nu$ as in
Table~\ref{tab:constants}) and the exact Gaussian moments:
\[
  \frac{\mathrm{norm}(L)}{L}\le
  \frac{\sqrt{\pi/2}}{2\pi}-\frac{L_1}{4\pi^2}
  +\frac{\nu\sqrt{\pi/2}\,L_2^2}{8\pi^3} .
\]
The second is a left-rectangle sum in the variable $v=Tt$: on the $t$-panel
corresponding to $v\in[ih,(i{+}1)h]$ one has $p(t)\le p(L_1 ih/(2\pi))$ and
$e^{-T^2t^2/2}\le e^{-(ih)^2/2}$, and the panel width is at most $Lh/(2\pi)$,
giving
\[
  \frac{\mathrm{norm}(L)}{L}\le\frac{h}{2\pi}
  \sum_{i<m}p\Bigl(\frac{L_1\,ih}{2\pi}\Bigr)e^{-(ih)^2/2} .
\]
The sum is evaluated in interval arithmetic by a recursion that carries
$e^{-(ih)^2/2}=x^{i^2}$ with $x=e^{-h^2/2}$, so that only one exponential is ever
computed.
\formal{nb\_step, nb\_loop, nb\_rect, nb\_closed, normalUpper, normalUpper\_sound, normalIntegral\_le, gaussMoment\_*}{Upper0423/Uniform/NormalBound.lean, Upper0423/LargeN/Normal.lean}

\paragraph{The low coefficients.} Substituting $v=2\pi t/L$ turns
$\mathrm{lowG}$ into $\int_0^{\pi/(2L)}
s|K(s)|\,\tfrac{v^2}{12}\,(E_1+E_2)\dd v$ with $s=Lv/(2\pi)$, and similarly for
$\mathrm{lowK}$. On each $v$-panel $[\alpha,\beta']\subseteq[0,V]$:
\begin{itemize}[leftmargin=1.6em]
\item the kernel factor $s|K(s)|$ is bounded by a prefix maximum over a
  precomputed table on the grid $s=j/(4N)$, $N=1024$; each table entry uses
  \eqref{eq:K-left} in the form
  $s^2|K(s)|^2=\sigma^2+(\tfrac1\pi-\sigma q(\pi s))^2$ with
  $\sigma=s(1-s)$, both factors being monotone on the relevant range;
\item $E_1\le e^{-\varrho v^2/2}$ and
  $E_2\le e^{-\varrho v^2/2+\varrho\kap L_e v^3}$ with
  $\varrho=1-h_{\mathrm{cap}}^2$ and
  $L_e=\min\bigl(L_2,\pi/(2\alpha)\bigr)$, the second using
  $v\le\pi/(2L)$;
\item the resulting smooth majorant is integrated by the fourth-order panel rule
  of Lemma~\ref{lem:panel}.
\end{itemize}
When $V=8<\pi/(2L)$ a closed-form Gaussian tail is added beyond $8$.
\formal{lb\_table, lb\_panel, lb\_tail, lowUpper, lb\_panel\_sound, lb\_tail\_sound, lowUpper\_sound}{Upper0423/Uniform/LowBound.lean}

\subsection{The cells}

\begin{theorem}
  \label{thm:uniform}
  Let $r=16$, $m=256$. There is a list of $282$ rational edges
  \[
    \tfrac1{100}=e_0<e_1<\dots<e_{281}=\tfrac{260889}{100000}
  \]
  such that every consecutive cell $[e_i,e_{i+1}]$ is accepted, and such that the
  initial cell $(0,\tfrac1{100}]$ is accepted in the variant form in which
  $(\mathrm{high}+\mathrm{tail})/L\le1/1000$ is proved analytically. Consequently,
  for every admissible $\mu$, every $n\ge256$ with
  $\beta/\sqrt n\le\tfrac{13003}{10000}$, and every $x$,
  \[
    \mathcal E(\mu,n,x)\le\frac{423}{1000} .
  \]
\end{theorem}
\formal{cellAccepted, cellsAccepted, smallCellAccepted, bud\_cover, normalizedError\_le\_of\_cellsAccepted}{Upper0423/Uniform/Budget.lean}
\formal{chunk0 \dots chunk8, acc0 \dots acc8, normalizedError\_uniform}{Upper0423/Uniform/Data/Chunk0.lean \dots Chunk8.lean, Uniform/Data/All.lean}

\begin{proof}[Proof of the deduction]
  Set $L=(\beta+\eta)h$ with $h=1/\sqrt n$. Since $n\ge256=r^2$ we have
  $r\le\sqrt n$, so $hr\le1$. The upper end: $\eta\le\beta$ (from $\eta\le1\le\beta$)
  gives $L\le2\beta/\sqrt n\le2\cdot\tfrac{13003}{10000}=\tfrac{26006}{10000}$,
  which is at most the last edge $e_{281}=2.60889$. If $L\le\tfrac1{100}$ the
  initial cell applies; otherwise $L$ lies in one of the $281$ cells, found by a
  linear scan (\Lean{bud\_cover}). Proposition~\ref{prop:reduction},
  the certified bounds $F,G,K$ for that cell, and
  Lemmas~\ref{lem:gammaconstraint}--\ref{lem:budget} give
  $\mathcal E(\mu,n,x)<\tfrac{423}{1000}$.
\end{proof}

The $281$ cells are declared in nine files \file{Uniform/Data/Chunk}$k$
(\Lean{Chunk0}--\Lean{Chunk7} with $35$ cells each, \Lean{Chunk8} with one),
each discharged by a single \Lean{native\_decide}; consecutive chunks share an
edge and are glued by \Lean{cellsAccepted\_append}. A tenth
\Lean{native\_decide}, in \file{Uniform/Data/All.lean}, discharges
\Lean{smallCellAccepted 16 256}. The edges are geometric with ratio
$\approx1.02$ (the generator was invoked as
\verb|gen_uniform.py 256 102/100 35|).

\begin{remark}
  The cell widths matter: the certified $F$, $G$, $K$ are constant on a cell but
  the true integrals vary, so a cell that is too wide is rejected, while too many
  cells make the \Lean{native\_decide} evaluation slow. The interesting statistic
  is the largest certified cell value, i.e.\ the left-hand side of the budget
  inequality maximized over the cells. Evaluating the same Lean functions
  \Lean{highTailUpper}, \Lean{normalUpper}, \Lean{lowUpper} and \Lean{bud\_val}
  on all $281$ cells gives
  \[
    \max_{1\le i\le 281}\ \bigl(\text{budget value of the cell }
      [e_{i-1},e_i]\bigr) = 0.414962\ldots,
  \]
  and the initial cell $(0,\tfrac1{100}]$ gives $0.412306\ldots$, so the
  certificate clears $423/1000$ with a margin of about $0.008$. (This is an
  auxiliary evaluation of the same definitions, reported for orientation; it is
  not part of the formal proof, which needs only the strict cellwise inequality
  that \Lean{native\_decide} establishes. The figure agrees with the
  approximate value $0.4150$ recorded in \file{docs/UPPER\_0423.md}.)
\end{remark}

\section{Trust boundary}
\label{sec:trust}

\subsection{Axioms}

Running
\begin{leancode}
#print axioms BerryEsseen.Upper0423.berryEsseenConstant_le_0423
\end{leancode}
reports
\begin{itemize}[leftmargin=1.6em]
\item the three standard axioms of Lean's type theory: \Lean{propext},
  \Lean{Classical.choice}, \Lean{Quot.sound}; and
\item $579$ auxiliary axioms generated by \Lean{native\_decide}, one per finite
  certificate.
\end{itemize}
We verified the count $579$ independently by enumerating the \Lean{native\_decide}
invocations in the module tree: $507$ in \file{Covers/Data/} (one per cover file:
$253$ \file{BE} files and $254$ \file{DW} files), $62$ in \file{Small/Data/} (one
per cover file), $9$ in \file{Uniform/Data/Chunk}$k$ and $1$ in
\file{Uniform/Data/All.lean}; $507+62+9+1=579$.

There is no \Lean{sorry} anywhere in the development and no user-declared
axiom. The consequence of the $579$ auxiliary axioms is precise and limited: for
each of the $579$ closed Boolean expressions, the proof asserts that the
expression evaluates to \Lean{true}, on the evidence of Lean's compiler and
runtime rather than of Lean's kernel. The mathematical content --- the smoothing
inequality, all the characteristic-function estimates, the semantics of the
interval arithmetic and the jets, the panel rule, and every soundness theorem
relating a Boolean acceptance predicate to a real-analytic statement --- is
kernel-checked.

\subsection{What is and is not claimed}

\begin{itemize}[leftmargin=1.6em]
\item We claim that the Lean development in \file{BerryEsseen/Upper0423/} builds
  and that its final theorem is the statement displayed in
  Section~\ref{sec:notation}, with the axiom profile just described.
\item We claim no independent verification. In particular, the repository's
  trusted-comparator configuration permits only the three standard axioms and
  targets the separate lower bound $0.40\le C_{\mathrm{BE}}$; we have not run any
  comparator on the $0.423$ theorem, and a configuration allowing only the three
  standard axioms would reject it.
\item We claim no peer review of this document or of the development.
\item We make no claim about the correctness of \cite{XiaoLi2026} or of
  \cite{IndependentBE2026}.
\item The definition of the constant being bounded is the one displayed in
  Section~\ref{sec:notation}. Readers who wish to audit only the statement need
  read only \file{Challenge.lean}, which imports Mathlib alone.
\end{itemize}

\subsection{Residual risks}

The honest list of things that could still be wrong is:
\begin{enumerate}[label=(\arabic*),leftmargin=2.2em]
\item a bug in Lean's compiler or runtime affecting the evaluation of one of the
  $579$ Boolean expressions (mitigated by the fact that the expressions use only
  arbitrary-precision integer and rational arithmetic, lists, arrays and
  structural recursion);
\item a bug in Lean's kernel or in Mathlib (the standard risk of any
  formalization);
\item a mismatch between the definition of \Lean{berryEsseenConstant} and the
  informal statement of Theorem~\ref{thm:main} (mitigated by the verbatim
  reproduction in Section~\ref{sec:notation});
\item an error in this document's \emph{description} of the Lean proof. The Lean
  sources, not this document, are authoritative; every lemma above carries the
  name and file of its formal counterpart so that the two can be compared.
\end{enumerate}

\section{Reproduction}
\label{sec:repro}

The development is built with
\begin{leancode}
lake build BerryEsseen.Upper0423.Main
\end{leancode}
The default library target \file{BerryEsseen} imports it, so
\texttt{lake build BerryEsseen} builds it as well. A full build is dominated by
the evaluation of the $579$ certificates.

The Lean data files are \emph{generated}, by the scripts in
\file{scripts/upper\_0423/}:
\begin{itemize}[leftmargin=1.6em]
\item \file{run\_covers.sh} (with \file{make\_jobs.py}, \file{cover\_gen.py})
  searches the covers using the research producer
  \file{research/native\_cover.py} and its C\texttt{++} engine;
\item \file{gen\_covers.py}, \file{gen\_all.py} write \file{Covers/Data/*.lean}
  and \file{Covers/All.lean};
\item \file{gen\_small\_cover.py}, \file{gen\_small\_data.py} write
  \file{Small/Data/*.lean} from the research covers;
\item \file{gen\_uniform.py 256 102/100 35} writes \file{Uniform/Data/*.lean}.
\end{itemize}
The search is entirely untrusted: it chooses the tree shapes, the smoothing
parameters $c$ and $\tau$, the panel subdivisions and the branch choices, but
every resulting numerical bound is recomputed from scratch by the Lean checkers,
whose soundness is proved for all accepted inputs. In particular a faulty search
can only produce a rejected certificate, never a false theorem.

\section{Related and concurrent work}
\label{sec:related}

\paragraph{Xiao--Li.} The repository \cite{XiaoLi2026} states the bound
$C_{\mathrm{BE}}\le 879/2000=0.4395$ with a Lean~4 formalization, also using
\Lean{native\_decide} for its certificates. Its analytic framework is a direct
ancestor of the present one; see Section~\ref{sec:attribution} for the itemized
attribution. As stated there, neither its numerical result nor its certificates
nor its Lean code are used in Theorem~\ref{thm:main}, and we make no assessment
of its correctness. It is an unpublished repository, with an accompanying note,
not a peer-reviewed publication.

\paragraph{The independent-summand constant.} The repository
\cite{IndependentBE2026} treats the analogous constant $C_{\mathrm{ind}}$ for
independent, not necessarily identically distributed, summands, and states
$C_{\mathrm{E}}\le C_{\mathrm{ind}}<0.44995$ along with several related results.
This is a different constant (necessarily $C_{\mathrm{ind}}\ge C_{\mathrm{BE}}$),
and nothing from that collection is used here. It is likewise an unpublished
collection of preprints, and we make no assessment of its correctness.

\paragraph{Relation to the published record.} Against
Table~\ref{tab:history}, Theorem~\ref{thm:main} would improve the best published
bound $0.4690$ of \cite{Shevtsova2013}. We stress again that we neither use nor
reprove any of the entries of that table, and that the present document is not
peer reviewed.

\paragraph{Relation to an earlier formalization.} The repository previously
contained a formalization of the weaker bound $C_{\mathrm{BE}}\le0.4688$ that
avoided \Lean{native\_decide}. That upper-bound proof has been removed; it remains
available in the repository's Git history. The $0.423$ development reuses
general-purpose modules originally written for it (the definitions, the Prawitz
smoothing inequality and its integrability, kernel estimates, characteristic-function
anchors and the Cantelli bound), which remain in the repository together with the
formalized lower bound $0.40\le C_{\mathrm{BE}}$.

\section*{Acknowledgements}

The analytic framework builds on the work of Haonan Xiao and Chunlin Li
\cite{XiaoLi2026} as detailed in Section~\ref{sec:attribution}. The
formalization rests entirely on Mathlib \cite{Mathlib}. The historical data of
Table~\ref{tab:history} follow the survey entry \cite{Tao2026}.

\bibliographystyle{plain}
\bibliography{references}

\end{document}
